\documentclass{article}
\usepackage{iclr2027_conference,times}

\usepackage{url}
\usepackage{multirow}
\usepackage{graphicx}
\usepackage{amssymb}
\usepackage{amsmath}
\usepackage{amsthm}
\usepackage{booktabs}
\usepackage{tabularx}
\usepackage{caption}
\usepackage{subcaption}
\usepackage{xcolor,colortbl}
\usepackage{array}
\usepackage{placeins}
\usepackage{longtable}
\usepackage{tikz}
\usetikzlibrary{arrows.meta,positioning,calc}
\usepackage{pgfplots}
\pgfplotsset{compat=1.16}
\usepackage[hidelinks]{hyperref}

\definecolor{cadd}{HTML}{E69F00}   % additive baseline
\definecolor{cclamp}{HTML}{0072B2} % clamp
\definecolor{cot}{HTML}{009E73}    % optimal transport
\definecolor{cgate}{HTML}{CC79A7}  % gated / NP-BW
\definecolor{ccal}{HTML}{009E73}
\pgfplotsset{
  qpanel/.style={width=.49\linewidth, height=3.4cm, grid=major,
    grid style={dashed,gray!14}, xmin=-35, xmax=95, ymin=-55, ymax=38,
    tick label style={font=\tiny},
    every axis label/.append style={font=\scriptsize},
    title style={font=\scriptsize\bfseries, yshift=-2pt}},
  bgpt/.style={only marks, mark=*, mark size=0.6pt, gray!30},
  stillocw/.style={only marks, mark=*, mark size=1.5pt, cadd},
  stillcr/.style={only marks, mark=*, mark size=1.5pt, cclamp},
  arocw/.style={-{Stealth[length=4pt]}, cadd!88!black, line width=1pt,
                quiver={u=\thisrow{u}, v=\thisrow{v}}},
  arcr/.style={-{Stealth[length=4pt]}, cclamp!88!black, line width=1pt,
               quiver={u=\thisrow{u}, v=\thisrow{v}}}
}
\newcommand{\qdecor}{%
  \fill[cgate!11] (axis cs:-35,-10.08) rectangle (axis cs:95,38);
  \draw[cgate!65!black, dashed, thick] (axis cs:-35,-10.08) -- (axis cs:95,-10.08);
  \node[anchor=north east, font=\tiny\itshape, text=cgate!60!black] at (axis cs:93,36) {gate ON};
  \node[anchor=south east, font=\tiny\itshape, text=gray] at (axis cs:93,-53) {gate OFF};
  \draw[ccal!75!black, thick, dash dot, rotate around={150:(axis cs:31,-4.6)}]
    (axis cs:31,-4.6) ellipse [x radius=30.5, y radius=13.7];
  \node[font=\tiny\itshape, text=ccal!60!black, anchor=west] at (axis cs:-30,-30) {calibrated zone};
  \addplot[bgpt] table[x=x,y=y,col sep=comma]{pic/viz/plane_mmlu.csv};
}
\pgfplotsset{
  planepanel/.style={width=.52\linewidth, height=4.6cm, grid=major,
    grid style={dashed,gray!15}, xmin=-35, xmax=95, ymin=-65, ymax=45,
    tick label style={font=\tiny}, every axis label/.append style={font=\scriptsize},
    title style={font=\scriptsize\bfseries, yshift=-2pt}},
  ptocw/.style={only marks, mark=*, mark size=1.1pt, cadd},
  ptcr/.style={only marks, mark=*, mark size=1.1pt, cclamp},
  ptrest/.style={only marks, mark=*, mark size=0.7pt, gray!55},
  disp/.style={-stealth, cgate, thick, quiver={u=\thisrow{u}, v=\thisrow{v}}}
}
\makeatletter\newcommand\inputtab[1]{\@@input #1.tex }\makeatother
\newtheorem{theorem}{Theorem}
\newtheorem{corollary}{Corollary}
\newtheorem{proposition}{Proposition}
\newtheorem{definition}{Definition}
\newtheorem{remark}{Remark}
\newcommand{\ocw}{OCW}
\newcommand{\crok}{CR}
\newcommand{\Wtwo}{W_2}
\captionsetup{font=small}
\setlength{\textfloatsep}{8pt plus 2pt minus 3pt}
\setlength{\intextsep}{8pt plus 2pt minus 3pt}
\setlength{\abovecaptionskip}{5pt}
\emergencystretch=1.8em
\renewcommand{\UrlFont}{\ttfamily\small}

\title{Projection-Transport Steering:\\ Conditional Control of Reasoning Behavior in LLMs}
\author{Anonymous authors\\Paper under double-blind review}

\begin{document}
\maketitle

\begin{abstract}
Activation steering changes a language model's behavior at inference time by
editing its residual stream, but most edits in use act on every input. An edit
strong enough to remove a behavior where it occurs therefore also changes the
inputs that never showed it. We propose Projection-Transport Steering (PTS), a
framework that separates when to intervene from how. The action is the
optimal-transport map from the projection onto a low-dimensional behavioral
subspace to a target law, which we show is the unique minimum-displacement edit
that leaves the rest of the residual stream untouched. The decision is a
likelihood-ratio test on a learned detection statistic. We show
that the selectivity of any conditional edit decomposes exactly into the
decision's error rates and the action's conversion rates, so one ungated run
predicts the selectivity of its gated configurations. On overconfident
multiple-choice reasoning with Gemma-2-2B, PTS is the most selective of eight
methods on held-out questions and lowers calibration error. On free-text toxicity in two model
families, PTS decides from the prompt before generating and is more selective
than the global shift, prompting, and refusal on the same decision.
\end{abstract}

\section{Introduction}
\label{sec:intro}

Many high-level behaviors of large language models appear to be carried by
directions in their activations. Activation steering
methods~\citep{turner2023actadd,rimsky2024caa,zou2023repe,li2023iti} read such a
direction off contrastive examples and add it to the residual stream at
inference time, at the cost of one vector addition per token and no
retraining. However, these methods apply the same edit to every input, so an
edit strong enough to remove a behavior also acts on the inputs that never
showed it: a refusal vector at a useful strength raises refusal on harmful and
harmless prompts alike~\citep{arditi2024refusal}.

The problem is sharper for a behavior that the model shows only on some inputs,
such as overconfidence in multiple-choice reasoning. The answers to change and
the answers to protect are both confident, one wrong and one right, and on
\texttt{gemma-2-2b-it} the global shift that eliminates every
overconfident-wrong answer also eliminates every confident-right one
(Figure~\ref{fig:frontier}); on the tuning split it costs fifteen accuracy points. A smaller coefficient
trades one failure for the other, because a shift moves the bulk of the
distribution as far as it moves the tail.

In this paper we introduce Projection-Transport Steering (PTS), a framework
that separates when to intervene from how. The \emph{action} is a transfer
function on the projection of the residual stream onto a behavioral subspace,
i.e., a few directions read off the model's own labeled answers. The
\emph{decision} is a likelihood-ratio test on a detection statistic fitted to
tell the two confident states apart, and the action runs only where the test
fires. Table~\ref{tab:family} places published methods in this frame: they
differ in the transfer function and in the subspace it acts on, and of those
listed only CAST and PCHI decide when to act.

Given the subspace and a target law, i.e., the distribution the steered
projection should follow, the action is determined: the optimal-transport map
from the unsteered law of the projection to the target law is the unique
minimum-displacement edit that leaves the orthogonal complement untouched (Theorem~\ref{thm:ot}). Additive steering is
this map exactly when the target is a translate of the unsteered law, and
ablation and clamping are the same map for other targets.

We score an edit by its selectivity, the share of overconfident-wrong answers
it removes minus the share of confident-right answers it destroys. The
selectivity of any conditional edit equals the decision's true-positive rate
times the rate at which the action converts flagged target answers, minus its
false-positive rate times the rate at which it converts flagged protected ones
(Proposition~\ref{prop:account}). With conversion rates from one ungated run,
this identity predicts the selectivity of gated configurations. When the
action converts both states at the same rate, the separability of the
detection statistic on unsteered answers caps the selectivity
(Theorem~\ref{thm:gate}), and the same quantities determine when a steering
action beats overriding every flagged output (Eq.~\eqref{eq:override}).

A decision that scores a finished answer has to regenerate every input it
flags. The same test can read the prompt or a short prefix, so PTS decides
before the answer and generates each input once, at two dot products and one table lookup
per token (Proposition~\ref{prop:seq}). Inputs it does not flag are generated
without any edit; on HumanEval the prompt decision flags no problem, and pass@1
is exactly its unsteered value (Appendix~\ref{app:capability}).

Our main contributions are:
\begin{itemize}\setlength{\itemsep}{1pt}
\item A framework and its theory: the optimal-transport characterization of
projection-local steering and an exact decomposition of conditional
selectivity (Sec.~\ref{sec:method}), which predicts measured gated selectivity
on both model families (Sec.~\ref{ssec:law}).
\item Held-out evidence on overconfident reasoning (Sec.~\ref{sec:results}):
with every configuration committed before the held-out run, PTS is the most
selective of eight methods on 3{,}000 MMLU questions of \texttt{gemma-2-2b-it}
and lowers the share of overconfident-wrong answers and the calibration error
at non-inferior accuracy. On \texttt{Qwen2.5-7B-Instruct} and on ARC-Challenge
it is level with or above every baseline after Holm adjustment, and in the
setting of PCHI it is more selective than PCHI under the same probe
(Sec.~\ref{sec:transfer}).
\item Free-text control decided from the prompt (Sec.~\ref{sec:transfer}): on
toxic continuation in two model families, PTS is more selective than the global
shift, prompting, and refusal on the same decision, and than CAST on Qwen, at
$1.09\times$ the unsteered wall clock on Gemma.
\end{itemize}

\begin{table}[t]\footnotesize\setlength{\tabcolsep}{5pt}
\centering
\resizebox{\linewidth}{!}{%
\begin{tabular}{l|l|c|c}
\hline
method & transfer function $g(p)$ & decision & cost / token \\
\hline
ActAdd, CAA~\citep{turner2023actadd,rimsky2024caa} & $p+\alpha$ & none & $\Theta(d)$ \\
ITI~\citep{li2023iti} & $p+\alpha$ on probe-selected heads & none & $\Theta(d)$ \\
RepE~\citep{zou2023repe} & $p+\alpha(x)$ & none & $\Theta(d)$ \\
directional ablation~\citep{arditi2024refusal} & $0$ & none & $\Theta(d)$ \\
SAE clamping~\citep{templeton2024scaling} & $c$ & none & $\Theta(d\,d_{\mathrm{sae}})$ \\
MiMiC~\citep{singh2024mimic} & full-space affine map & none & $\Theta(d^2)$ \\
Linear-AcT~\citep{rodriguez2025act} & $(1{-}\lambda)p+\lambda(\omega p+\beta)$, per neuron & none & $\Theta(d)$ \\
CAST~\citep{lee2024cast} & $p+\alpha$ & prompt similarity & $\Theta(d)$ \\
PCHI~\citep{liu2026probe} & learned head rescaling & probe at the readout & $\Theta(d)$ \\
\hline
\textbf{PTS (ours)} & $S(p)$, the transport map to $\nu$ & likelihood-ratio test & $\Theta(kd)$ \\
\hline
\end{tabular}}
\caption{\textbf{Most published steering operators apply one transfer function
to every input.} Each row applies a transfer function $g$ to a $k$-dimensional
projection $p$ of the residual stream, with $k{=}1$ on a behavioral direction
for most rows, $k{=}d_{\mathrm{sae}}$ for SAE clamping and $k{=}d$ for the
full-space maps. The rows differ in the shape of $g$ and in the $k$ it acts on,
and only CAST and PCHI decide when to apply it. Theorem~\ref{thm:ot} identifies
the cost-minimal $g$ for a target law $\nu$, and
Proposition~\ref{prop:account} writes the selectivity of every row with a
decision in terms of that decision and its action.}
\label{tab:family}
\end{table}

\begin{figure}[t]
\centering
\begin{tikzpicture}
\begin{axis}[width=.92\linewidth, height=6.0cm,
  xlabel={removal: share of overconfident-wrong answers removed},
  ylabel={retention: confident-right kept},
  xmin=-0.02, xmax=1.03, ymin=-0.03, ymax=1.03, grid=major, grid style={dashed,gray!18},
  tick label style={font=\scriptsize}, every axis label/.append style={font=\footnotesize},
  legend columns=2, legend style={font=\scriptsize, draw=none, fill=white, fill opacity=.9, text opacity=1,
    at={(0.01,0.01)}, anchor=south west, column sep=4pt}]
\foreach \s in {0,0.1,0.2,0.3} {\addplot[gray!55, densely dotted, forget plot, domain=\s:1] {1+\s-x};}
\node[font=\tiny, gray, rotate=-38] at (axis cs:0.83,0.47) {Sel $=0.3$};
\node[font=\tiny, gray, rotate=-38] at (axis cs:0.64,0.36) {Sel $=0$};
\addplot[only marks, mark=*, mark size=0.9pt, gray!55] table[x=removal, y=retention, col sep=comma]{pic/frontier_v4_confirm_random.csv};
\addlegendentry{same decision, random direction ($\times100$)}
\addplot[cadd, thick, mark=*, mark size=1.4pt] coordinates {(0,1) (0.5587,0.6612) (0.8411,0.4153) (0.9813,0.1279) (1.0,0.0)};
\addlegendentry{global shift, increasing dose}
\addplot[only marks, mark=square*, mark size=1.7pt, cclamp] table[x=removal, y=retention, col sep=comma]{pic/frontier_v4_confirm_baseline.csv};
\addlegendentry{prompting, ablation, MiMiC, Linear-AcT}
\addplot[only marks, mark=triangle*, mark size=2.3pt, cot] table[x=removal, y=retention, col sep=comma]{pic/frontier_v4_confirm_cast.csv};
\addlegendentry{CAST, prompt and trace condition}
\addplot[only marks, mark=diamond*, mark size=2.4pt, cgate!80!black] table[x=removal, y=retention, col sep=comma]{pic/frontier_v4_confirm_pts_onepass.csv};
\addlegendentry{PTS, decision before the answer}
\addplot[only marks, mark=star, mark size=4pt, very thick, cgate!70!black] table[x=removal, y=retention, col sep=comma]{pic/frontier_v4_confirm_pts.csv};
\addlegendentry{PTS, post-hoc}
\end{axis}
\end{tikzpicture}
\caption{\textbf{Gating the action by a decision places PTS above the dose
curve of the global shift.} Every method is shown on the 3{,}000 held-out MMLU
questions of \texttt{gemma-2-2b-it}, at the configuration it won on a separate
tuning split. The orange curve is the global shift at
increasing dose: it gains removal only by losing confident-right answers, and
its last point removes every overconfident-wrong answer and keeps none of the
confident-right ones. Gray points keep the PTS decision and dose and replace
its direction by one of a hundred random unit directions; post-hoc PTS is more
selective than all of them. Dotted lines are levels of selectivity, removal
minus the share of confident-right answers lost.}
\label{fig:frontier}
\end{figure}


\section{Related work}
\label{sec:related}

\textbf{Activation steering.} ActAdd~\citep{turner2023actadd} adds the
activation difference of one contrast prompt pair to the residual stream, and
CAA~\citep{rimsky2024caa} averages it over many pairs.
ITI~\citep{li2023iti} shifts attention heads selected by linear probes,
RepE~\citep{zou2023repe} reads and writes representation directions,
directional ablation~\citep{arditi2024refusal} removes one direction from every
activation, and SAE clamping~\citep{templeton2024scaling} fixes one sparse
feature. AxBench~\citep{wu2025axbench} finds prompting stronger than
such edits, and reliability
studies~\citep{tan2024generalisation,braun2025unreliability} find that their
effect varies widely across inputs. All of them apply one edit to every input,
with a shape fixed by hand; Theorem~\ref{thm:ot} derives the shape from a
target law.

\textbf{Transport and distribution matching.} MiMiC~\citep{singh2024mimic}
derives the affine map that matches the mean and covariance of two activation
distributions, the Gaussian optimal-transport map, and applies it to the full
activation. Linear-AcT~\citep{rodriguez2025act} fits a univariate affine
transport map to each neuron with a strength $\lambda\in[0,1]$ and writes
earlier edits as special cases; LinEAS~\citep{rodriguez2025lineas} learns such
maps across layers end to end. COAST~\citep{nguyen2026coast} moves each
activation, at fixed norm, to a fixed cosine with a target direction at the
smallest second-moment-weighted displacement, which in our terms transports
that cosine onto a point mass, and A-LQR~\citep{nguyen2026alqr} casts steering
as linear-quadratic control. Unlike these maps, PTS transports only a
$k$-dimensional projection and leaves its orthogonal complement fixed, which
reduces the cost-minimal edit, for any target law, atoms included, to
optimal transport in $k$ dimensions, in closed form for $k=1$.

\textbf{Deciding when to act.} CAST~\citep{lee2024cast} compares the prompt's
hidden state with a condition vector at a grid-searched threshold and adds a
behavior vector when the comparison holds; we fit its condition on the labeled
pool of our probe. DSAS~\citep{ferrando2026dsas} decouples when to
steer from how: a logistic gate scales any steering map per token.
SADI~\citep{wang2025sadi}, dynamic activation
composition~\citep{scalena2024dynamic} and GAPS~\citep{gaps2026} adapt the edit
per input, per generation step and per dimension, respectively.
PCHI~\citep{liu2026probe} conditions a learned rescaling of downstream
attention heads on a frozen probe read at a confidence-template token; we run
PTS in its setting (Sec.~\ref{sec:transfer}). Activations predict whether an
answer is right~\citep{kadavath2022know,azaria2023lying} and the first
generated tokens whether steering will succeed~\citep{steerable2026}; probes
of a finished behavior forecast it poorly from a
prefix~\citep{fpcg2026}, and a steering write earlier in the same pass can
degrade a direction read after it~\citep{luo2026yopo}. Closer to our work,
CAST, DSAS and PCHI already separate the decision from the action; PTS adds an
exact account of the pair (Prop.~\ref{prop:account}), with the ceiling and the
reference arms it implies, which we run for every conditional method.

\section{Method: Projection-Transport Steering}
\label{sec:method}

A conditional steering operator has to settle what edit to apply and to which
inputs. PTS answers with a transport map on a behavioral subspace and a
likelihood-ratio test on a detection statistic, and accounts for how the two
compose. Proofs are in Appendix~\ref{app:proofs}.

\textbf{The action.} Let $h\in\mathbb{R}^d$ be the residual-stream activation
at a fixed decoder layer, and let $V\in\mathbb{R}^{k\times d}$ have orthonormal
rows spanning what we call the behavioral subspace; $q=Vh$ is the projection.
In our experiments $k=1$, and the row of $V$ is the unit difference of mean
activations between confident and hesitant answers (toxic and clean
continuations on free text), which we call the confidence coordinate. We write
$\mu$ for the law of $h$ under unsteered generation, $\mu_V$ for the induced
law of $q$, and $\nu$ for a target law on $\mathbb{R}^k$. Our goal is an edit
that gives the projection the law $\nu$, leaves the orthogonal complement
untouched, and moves $h$ as little as possible, since large edits are what
damage generation (Sec.~\ref{sec:results}).

\begin{definition}[admissible steering]
\label{def:local}
$T:\mathbb{R}^d\to\mathbb{R}^d$ is admissible for $(V,\nu)$ if
(i) $(I-V^\top V)\,T(h)=(I-V^\top V)\,h$ $\mu$-a.s.\ (\emph{locality}: the
orthogonal complement is untouched), and (ii) $(V\!\circ T)_{\#}\mu=\nu$
(\emph{goal}: steered projections follow the target law). Its cost is
$C(T)=\mathbb{E}_\mu\lVert T(h)-h\rVert^2$.
\end{definition}

\begin{theorem}[steering is optimal transport of the projection]
\label{thm:ot}
Let $\mu_V$ be absolutely continuous with finite second moment and let $\nu$
have a finite second moment. Then $\min_T C(T)=\Wtwo^2(\mu_V,\nu)$ over
admissible $T$, attained by
\begin{equation}
T^*(h)=h+V^\top\big(S(Vh)-Vh\big),
\label{eq:edit}
\end{equation}
where $S$ is the ($\mu_V$-a.s.\ unique) Brenier map from $\mu_V$ to $\nu$.
Every optimal $T$ satisfies $VT(h)=S(Vh)$, hence $T=T^*$ by locality,
$\mu$-almost surely.
\end{theorem}

In words, once the subspace and the target law are fixed, the edit is
determined: it moves only the projection, by the transfer function $S$ ($g$ in
Table~\ref{tab:family}). For $k=1$, $S=F_\nu^{-1}\circ F_{\mu_V}$ is the
monotone rearrangement, stored as a 41-knot piecewise-linear lookup;
atomlessness of $\mu_V$ then suffices and $\nu$ may have atoms. Directional ablation is the map onto $\delta_0$, a clamp at $c$ the
map onto the law of $\min(q,c)$, and the shift $S(q)=q+\alpha$ the map
onto a translate of $\mu_V$; for Gaussian marginals $S$ is
affine~\citep{knott1984optimal,gelbrich1990formula}. Since every nondecreasing
transfer function is the optimal map onto its own image law,
Theorem~\ref{thm:ot} does not rank these edits: it turns the choice of edit
into a choice of target law, made on a tuning split (Sec.~\ref{sec:exp}), and
it concerns displacement alone. Read backwards, it says when the shift is
optimal.

\begin{corollary}[when additive steering is optimal, and its limit]
\label{cor:additive}
(a) $S(q)=q+\alpha$ is the optimal map iff $\nu$ is the translate of $\mu_V$
by $\alpha$.
(b) Let two subpopulations have the same law of $q$. If, under an edit, the
probability that an input leaves its state depends on $h$ only through $q$,
the edit changes the state of both at the same rate.
\end{corollary}

\textbf{The decision and its accounting.} The states we must tell apart,
overconfident-wrong (\ocw{}) and confident-right (\crok{}), are both confident,
and their laws on the confidence coordinate nearly coincide
(Sec.~\ref{sec:results}), so Corollary~\ref{cor:additive}(b) predicts that any
transfer function of that coordinate changes both at similar rates, as the
global shift does (Fig.~\ref{fig:frontier}). A decision therefore reads a second statistic $s$ of the input and runs the action on an
event $\Gamma$ determined by $s$, in practice $\{s>\tau\}$ for a scalar score;
off $\Gamma$ the output is exactly the unsteered one. We measure the result with selectivity,
$\mathrm{Sel}=P(\text{leave \ocw}\mid\text{\ocw})-P(\text{leave \crok}\mid\text{\crok})$, the
share of \ocw{} answers removed minus the share of \crok{} answers lost; on
free text, toxic and clean continuations take their places. Write
$\mathrm{TPR}$ and $\mathrm{FPR}$ for the rates at which the decision fires on
\ocw{} and \crok{}, and $\rho_O,\rho_C$ for the rates at which the action
changes the state of a flagged \ocw{} and a flagged \crok{} answer.

\begin{proposition}[selectivity decomposes]
\label{prop:account}
For any decision and action that leave unflagged outputs unchanged,
\begin{equation}
\mathrm{Sel}=\mathrm{TPR}\cdot\rho_O-\mathrm{FPR}\cdot\rho_C .
\label{eq:account}
\end{equation}
\end{proposition}

The proof is one line of total probability. The decision sets
$(\mathrm{TPR},\mathrm{FPR})$, and the action sets $(\rho_O,\rho_C)$ on the
flagged answers. Since a post-hoc gated
operator returns, input by input, the unsteered or the ungated output, one
ungated run per action composes every such configuration, and the ungated
rates predict its selectivity when the action converts flagged and unflagged answers of a
state alike (Sec.~\ref{ssec:law}). The decomposition also fixes three reference
arms on the same decision: a no-edit arm measures what regeneration alone
changes, a random direction at the same dose what any perturbation of that
size changes, and an override of the reported confidence of every flagged
answer ($\rho_O=\rho_C=1$; on free text, a refusal) has selectivity
$\mathrm{TPR}-\mathrm{FPR}$. The action beats that
override exactly when
\begin{equation}
\mathrm{TPR}\,(1-\rho_O)<\mathrm{FPR}\,(1-\rho_C),
\label{eq:override}
\end{equation}
that is, when the share of protected answers flagged but spared exceeds the
share of target answers flagged but not converted.

\begin{theorem}[optimal decision and selectivity ceiling]
\label{thm:gate}
Let $p_O,p_C$ be the densities of $s$ on \ocw{} and \crok{} answers, and
$r_O(s),r_C(s)$ the probabilities that the action converts such answers at
$s$.
(a) If $r_O$ and $r_C$ do not depend on $s$, then among decisions with equal
$\mathrm{FPR}$ the likelihood-ratio test on $s$ maximizes $\mathrm{Sel}$
(Neyman--Pearson); under equal-covariance Gaussian
class conditionals it is the linear discriminant. In general, among decisions
with equal $\mathrm{FPR}\cdot\rho_C$, the maximizer thresholds
$p_O(s)\,r_O(s)/\big(p_C(s)\,r_C(s)\big)$.
(b) If $s$ is scalar and the concave hull of its ROC has area $\mathrm{AUC}$,
every decision $\{s>\tau\}$ has $\mathrm{TPR}-\mathrm{FPR}\le2\,\mathrm{AUC}-1$,
which caps the override on it and any action with $\rho_O=\rho_C$.
\end{theorem}

Part (b) is the ceiling of the introduction: computed from unsteered answers
alone, it bounds what a threshold decision contributes when the action
converts both states alike. The premise of (a) holds only approximately on our
data, where flagged answers convert more often than unflagged ones of the same
state (Appendix~\ref{app:law}); Appendix~\ref{app:proofs} shows that it
matters, so (a) motivates the estimator without certifying it. The decision is
a logistic-regression probe, an estimate of the log-likelihood ratio up to a
constant that $\tau$ absorbs, on the activations of the finished answer, its
first $t$ tokens or the prompt, fitted on labeled answers disjoint from every
tuning and evaluation split.

\begin{theorem}[transport under a fixed decision]
\label{thm:joint}
Let $P(\Gamma)>0$, let $\mu_V(\cdot\mid\Gamma)$ be absolutely continuous, and
let $\nu_\Gamma$ have a finite second moment. Among interventions that are
local to $V$, the identity off $\Gamma$, and whose projections on $\Gamma$
follow $\nu_\Gamma$, the cost-minimal one applies on $\Gamma$ the Brenier map
from $\mu_V(\cdot\mid\Gamma)$ to $\nu_\Gamma$, at cost
$P(\Gamma)\,\Wtwo^2(\mu_V(\cdot\mid\Gamma),\nu_\Gamma)$.
\end{theorem}

The action can thus be fitted to the flagged population without giving up
displacement optimality. The theorem holds the decision fixed: under a linear
gain and a price on displacement, a graded dose, the form of the per-token
interpolation of DSAS, beats the switch exactly when its optimum lies strictly
between $0$ and $1$ with positive probability (Prop.~\ref{prop:dose}). Scoring a finished answer and
regenerating the flagged ones generates those inputs twice; deciding from the
prompt or a prefix does not.

\begin{proposition}[single-pass decision]
\label{prop:seq}
Let $z_t$ be computed from the activations of the prompt, or of the prompt and
the first $t$ generated tokens. Among decisions measurable with respect to
$z_t$, Theorem~\ref{thm:gate} holds with $z_t$ in place of $s$, so the ceiling
of an early decision is read from the ROC of $z_t$ on unsteered answers. If the
action starts after the last position $z_t$ reads, the operator runs in one
forward pass: the decision reads only unedited activations, earlier positions
keep their cached keys and values, and for $k=1$ each token costs at most two extra
$d$-dimensional dot products, one lookup and one vector update.
\end{proposition}
\section{Experimental setup}
\label{sec:exp}\label{sec:data}

\textbf{Models and behavior.} The main model is
\texttt{gemma-2-2b-it}~\citep{gemma2024}. The confirmatory comparison is
repeated on \texttt{Qwen2.5-7B-Instruct}, a second model family on which
everything is refitted, and on ARC-Challenge~\citep{clark2018arc}, where
every decision is refitted with ARC training and validation questions added to
its pool. On MMLU~\citep{hendrycks2021mmlu} both models are prompted to reason
inside \texttt{<think>...</think>} and answer \verb|\boxed{letter}|, with
greedy decoding. The action runs at layer 14 of 26 (Gemma) and 14 of 28 (Qwen).
Section~\ref{sec:transfer} adds toxic continuation as a free-text behavior and
the setting of PCHI~\citep{liu2026probe}, the closest conditional method.

\textbf{Readout.} The primary readout (M5) scores the answer the model writes.
Correctness compares the boxed letter with the gold answer. Confidence is
$P(\mathrm{YES})$ when the model is then asked, in the same conversation,
whether that answer is correct, and a steered question is steered in both
turns. Correctness and whether $P(\mathrm{YES})\ge0.5$ define four states;
the two confident ones are overconfident-wrong (\ocw{}, wrong) and
confident-right (\crok{}, right). Appendix~\ref{app:secondary} reports two secondary
readouts, the letter logit of the same answer and a per-option verification,
and Appendix~\ref{app:prompts} gives the prompts.

\textbf{Splits.} The MMLU splits are disjoint subsets of the test set: 400
questions to extract the directions and transport targets, a pool of 4{,}000 to
fit the detection probe, 1{,}200 on which every configuration is tuned, and
3{,}000 used once for confirmation. ARC uses 1{,}012 probe and 400 tuning
questions from its training and validation sets and its 1{,}165 test questions
for confirmation.

\textbf{Metrics.} Each question is paired with itself before and after the
edit. Removal is $P(\text{not }\ocw\mid\ocw)$, retention is
$P(\crok\mid\crok)$, and $\mathrm{Sel}=\text{removal}-(1-\text{retention})$,
which is $1$ for an edit that removes every \ocw{} and keeps every \crok{} and
$0$ for one that treats the two states alike. Because selectivity conditions on
the starting state, we also report, over all questions, the change in the share
of \ocw{} answers, the 10-bin expected calibration error (ECE) of the stated
confidence, and the change in accuracy.

\textbf{Methods and tuning.} We compare prompting with a calibration
instruction, additive CAA~\citep{rimsky2024caa}, directional
ablation~\citep{arditi2024refusal}, MiMiC~\citep{singh2024mimic},
Linear-AcT~\citep{rodriguez2025act}, CAST with its prompt
condition~\citep{lee2024cast} fitted on the same labeled pool as our probe, and
a CAST-style trace condition, a shift conditioned on the finished trace that,
like post-hoc PTS, regenerates the questions it flags. All methods are reimplemented on our hooks
and scored with the same readout. PTS decides post-hoc, from the finished
unsteered answer, and regenerates the flagged questions; at the prompt, in one
forward pass over it before generation; or at a prefix of
$t\in\{16,32,64,128\}$ generated tokens, which keeps decision and edit in one
pass. Its action is a transport map of Theorem~\ref{thm:ot} on the confidence coordinate,
and its grids cover shifts at several doses and ablation. One rule selects
every configuration on the tuning split: the highest selectivity among
configurations whose accuracy change is at least $-0.01$, or the highest
overall when none meets that floor; every confirmatory number is a real run.
The budgets are one configuration for prompting and for ablation, eight for
CAA, nine each for MiMiC and Linear-AcT, and sixteen for each conditional
method, ours and the two CAST variants alike (Appendix~\ref{app:tuning}).

\textbf{Reference arms and statistics.} Every conditional arm is run against
three references that share its decision: a no-edit arm, which regenerates
the flagged inputs with the identity action; random unit directions at the same dose (ten
per conditional multiple-choice arm; one hundred for PTS, or thirty displaced
as far as the selected ablation on Qwen-MMLU); and an override that marks every
flagged answer as not confident, or replaces every flagged continuation by a
refusal on free text. Intervals are $95\%$ paired bootstrap intervals over
questions (prompts on free text) with $10^4$ resamples. The seven comparisons
of PTS with the baselines are two-sided and Holm-adjusted within each setting,
each later protocol Holm-adjusts its own tests, and an accuracy change is
non-inferior when the lower end of its interval is above $-0.02$. Every
protocol was committed before its held-out rollouts, the multiple-choice one
with the hashes of its selected configurations and fitted detectors
(Appendix~\ref{app:protocol}).

\section{Selective removal of overconfident answers}
\label{sec:results}

\begin{figure}[t]
\centering
\begin{tikzpicture}
\begin{axis}[name=law, width=.5\linewidth, height=4.8cm,
  xlabel={predicted from ungated rates}, ylabel={composed gated selectivity},
  xmin=-0.01, xmax=0.39, ymin=-0.01, ymax=0.39, grid=major, grid style={dashed,gray!20},
  tick label style={font=\scriptsize}, every axis label/.append style={font=\footnotesize},
  legend columns=2, legend style={font=\tiny, draw=none, fill=white, fill opacity=.85, text opacity=1, at={(0.01,0.99)}, anchor=north west, column sep=2pt},
  title={(a) selectivity decomposes}, title style={font=\footnotesize}]
\addplot[gray, dashed, forget plot] coordinates {(-0.01,-0.01) (0.39,0.39)};
\addplot[only marks, mark=*, mark size=0.8pt, cadd, opacity=.6] table[x=law, y=sim, col sep=comma]{pic/v4_law.csv};
\addlegendentry{shift}
\addplot[only marks, mark=square*, mark size=0.8pt, cclamp, opacity=.6] table[x=law, y=sim, col sep=comma]{pic/v4_law_ablate.csv};
\addlegendentry{ablation}
\addplot[only marks, mark=triangle*, mark size=1pt, cot, opacity=.6] table[x=law, y=sim, col sep=comma]{pic/v4_law_mimic.csv};
\addlegendentry{MiMiC}
\addplot[only marks, mark=diamond*, mark size=0.9pt, cgate, opacity=.6] table[x=law, y=sim, col sep=comma]{pic/v4_law_clamp.csv};
\addlegendentry{clamp}
\addplot[only marks, mark=o, mark size=0.9pt, gray, opacity=.6] table[x=law, y=sim, col sep=comma]{pic/v4_law_qwen.csv};
\addlegendentry{Qwen, all actions}
\addplot[only marks, mark=star, mark size=2.6pt, black, thick] table[x=law, y=real, col sep=comma]{pic/v4_law_real.csv};
\addlegendentry{confirmatory arms (real)}
\end{axis}
\begin{axis}[at={(law.east)}, anchor=west, xshift=1.2cm, width=.5\linewidth, height=4.8cm,
  xlabel={generated tokens read by the decision}, ylabel={held-out AUROC},
  xtick={0,32,64,128,256}, xticklabels={prompt,32,64,128,all}, ymin=0.6, ymax=0.82,
  grid=major, grid style={dashed,gray!20}, tick label style={font=\scriptsize},
  every axis label/.append style={font=\footnotesize},
  legend style={font=\tiny, draw=none, fill=none, at={(0.98,0.02)}, anchor=south east},
  title={(b) the decision can be taken early}, title style={font=\footnotesize}]
\addplot[cgate, thick, mark=*, mark size=1.2pt, error bars/.cd, y dir=both, y explicit]
  table[x=t, y=auroc, y error plus expr=\thisrow{hi}-\thisrow{auroc}, y error minus expr=\thisrow{auroc}-\thisrow{lo}, col sep=comma]{pic/v4_decision_gemma.csv};
\addlegendentry{Gemma-2-2B}
\addplot[cclamp, thick, dashed, mark=square*, mark size=1.2pt, error bars/.cd, y dir=both, y explicit]
  table[x=t, y=auroc, y error plus expr=\thisrow{hi}-\thisrow{auroc}, y error minus expr=\thisrow{auroc}-\thisrow{lo}, col sep=comma]{pic/v4_decision_qwen.csv};
\addlegendentry{Qwen2.5-7B}
\addplot[gray, densely dotted, thick, forget plot] coordinates {(0,0.669) (256,0.669)};
\node[font=\tiny, gray, anchor=south west] at (axis cs:2,0.669) {CAST-style trace condition};
\end{axis}
\end{tikzpicture}
\caption{\textbf{Gated selectivity is predicted from one ungated run, and the
decision can be read from the prompt.}
(a) Each point is a decision--action--threshold combination on the tuning split
(240 for Gemma in color, 240 for Qwen in gray). The horizontal axis is
$\mathrm{TPR}\,\rho_O-\mathrm{FPR}\,\rho_C$ with ungated conversion rates, the
vertical axis is the gated selectivity composed from the same ungated run, and
the dashed line is the identity. Stars are the four gated arms of the Gemma
confirmatory split, with their real gated runs on the vertical axis.
(b) Held-out AUROC between \ocw{} and \crok{} on the 3{,}000 confirmatory
questions of a probe that reads the prompt, the first $t$ generated tokens, or
the whole answer. Bars are $95\%$ bootstrap intervals, and the dotted line is
the AUROC of the trace direction used by the CAST-style trace condition.}
\label{fig:law}\label{fig:decision}
\end{figure}

On the Gemma tuning split the global shift at $-0.75\lVert
h\rVert$ removes every overconfident-wrong answer and every confident-right one
and lowers accuracy by $15.3$ points, and on the held-out questions the same
dose keeps none of the $1{,}399$ confident-right answers
(Fig.~\ref{fig:frontier}). Corollary~\ref{cor:additive}b suggests why: a
transfer function of one coordinate converts two states that share its law at
the same rate.

By Eq.~\eqref{eq:account} and Theorem~\ref{thm:gate}b, an override and any
action that converts both states at the same rate reach at most
$\max(\mathrm{TPR}-\mathrm{FPR})$ of the statistic they threshold, a ceiling computed
from unsteered answers before any steered rollout. On the confidence coordinate
the edit moves, the two confident states overlap heavily: the held-out ceiling
is $0.216$ on Gemma and $0.072$ on Qwen (AUROC $0.635$ and $0.53$). We read
this as the coordinate tracking how confident the model is, with little
information on whether the answer is right. The detection statistic of PTS, a
probe fitted on the separate pool of 4{,}000 labeled answers, raises the
ceiling to $0.370$ and $0.449$ (AUROC $0.741$ and $0.785$).

\subsection{Gated selectivity from one ungated run}
\label{ssec:law}

Proposition~\ref{prop:account} turns one ungated run into a prediction for
every decision placed in front of it. On the tuning split we measured the
conversion rates of six actions once, composed each with two decisions at 20
thresholds, and compared $\mathrm{TPR}\,\rho_O-\mathrm{FPR}\,\rho_C$ with the
selectivity composed from the same run (Fig.~\ref{fig:law}a). The two agree
with correlation $0.97$ over 240 combinations on Gemma and $0.91$ over 240 on
Qwen. A post-hoc gated arm is, question by question, either its unsteered or
its ungated output, so its selectivity can be composed from one ungated run up
to batched-decoding nondeterminism. Across twelve real gated arms on the three
held-out splits the composition lies within $0.025$ of the real run, with a
median gap of $0.005$ (Appendix~\ref{app:law}). We therefore score gated
configurations during tuning from one generation per action, and those that
decide at a prefix by real runs.

\subsection{PTS on held-out questions}
\label{ssec:confirm}

\begin{table}[t]\footnotesize\setlength{\tabcolsep}{3.5pt}
\centering
\resizebox{\linewidth}{!}{%
\begin{tabular}{l|ccc|cc|cc}
\hline
method & $\Delta$acc & removal & retention & Sel $\uparrow$ [95\% CI] & $\Delta$Sel ($p$) & $\Delta$\ocw{} (pp) $\downarrow$ & $\Delta$ECE $\downarrow$ \\
\hline
\inputtab{gen/confirm_gemma}
\hline
\end{tabular}}
\caption{\textbf{PTS is the most selective of eight methods on held-out
questions.} \texttt{gemma-2-2b-it}, 3{,}000 MMLU questions never used for
fitting or tuning ($n_{\ocw}{=}963$, $n_{\crok}{=}1{,}399$), each method at the
configuration selected for it on the tuning split under one rule. $\Delta$Sel
is the row minus PTS post-hoc (paired bootstrap, $10^4$ resamples), with
Holm-adjusted $p$ over the seven baselines; $\Delta$\ocw{} is the change in the
share of overconfident-wrong answers among all questions. The middle block
decides before the answer is written. The last two rows keep the PTS decision and
replace its action by an identity or by a random unit direction at the same
dose (mean over 100 directions). $^*$: $95\%$ interval excludes $0$;
$^\dagger$: the $95\%$ interval of $\Delta$acc reaches below $-0.02$.}
\label{tab:confirm}
\end{table}

In Table~\ref{tab:confirm}, post-hoc PTS reaches selectivity $0.294$ $[0.256,
0.334]$: it removes $49\%$ of overconfident-wrong answers and keeps $80\%$ of
confident-right ones. The closest baseline, the CAST-style trace condition,
reaches $0.248$, and the other six trail PTS by $0.074$ to
$0.233$; all seven differences remain significant after Holm adjustment. The
lead over it holds when the confidence cutoff moves from
$0.5$ to $0.4$ ($+0.044$ $[0.007, 0.080]$) or to $0.6$ ($+0.068$ $[0.028,
0.107]$).

Over all questions, PTS lowers the share of overconfident-wrong answers by
$15.0$ points $[13.7, 16.4]$ and the expected calibration error from $0.354$ to
$0.255$, with an accuracy change of $+0.003$ $[-0.009, 0.015]$, within the
non-inferiority margin of $-0.02$. Deciding before the answer gives $0.294$
from the first 32 generated tokens, in a single pass, and $0.306$ from the
prompt, the latter with an
accuracy interval that reaches the margin ($^\dagger$); the prompt probe
separates the two confident states nearly as well as the probe on the finished
answer (Fig.~\ref{fig:law}b).

\textbf{Reference arms.} Random unit directions sent through the PTS decision
at the PTS dose are less selective than the confidence direction. Over one
hundred such directions the mean is $0.131$, the 95th percentile $0.237$ and
the maximum $0.293$; PTS is above all of them and above their mean by $0.163$
$[0.140, 0.188]$ (gray points in Fig.~\ref{fig:frontier}). Regenerating the
flagged questions with no edit scores $0.049$, which measures what regeneration
in a new batch changes by itself. Appendix~\ref{app:refs} reports the override
arm for every conditional method.

\textbf{Capability away from the target.} Gated by the prompt decision, the
shift leaves HumanEval unchanged: the decision fires on none of the problems, so
pass@1 stays at its unsteered $0.396$, whereas the same shift applied to every
input at $-0.5\lVert h\rVert$, the dose selected for the prompt decision,
lowers it to $0.281$. Refitted with 600 AlpacaEval instructions as negatives,
the probes fire on none of the HumanEval and AlpacaEval prompts, so the operator
is the exact identity on code and open instructions, and with the selected
actions they reach held-out selectivity $0.280$ and $0.300$
(Appendix~\ref{app:capability}).
\section{PTS beyond the main setting}
\label{sec:transfer}

\begin{table}[t]\footnotesize\setlength{\tabcolsep}{3pt}
\centering
\resizebox{\linewidth}{!}{%
\begin{tabular}{l|cccc|cccc}
\hline
 & \multicolumn{4}{c|}{Qwen2.5-7B, MMLU ($n_{\ocw}{=}563$, $n_{\crok}{=}2{,}007$)} & \multicolumn{4}{c}{Gemma-2-2B, ARC ($n_{\ocw}{=}228$, $n_{\crok}{=}780$)} \\
method & $\Delta$acc & Sel $\uparrow$ [95\% CI] & $\Delta$Sel ($p$) & $\Delta$ECE $\downarrow$ & $\Delta$acc & Sel $\uparrow$ [95\% CI] & $\Delta$Sel ($p$) & $\Delta$ECE $\downarrow$ \\
\hline
\inputtab{gen/confirm_qa}
\hline
\end{tabular}}
\caption{\textbf{A second model family and a second benchmark.} Same protocol
and notation as Table~\ref{tab:confirm}. Qwen refits everything; ARC reuses
the MMLU directions and refits the decisions and configurations. The
random-direction row averages 100 directions on ARC and, because Qwen selects
ablation, 30 directions displaced as far as the ablation displaces the
confidence direction.}
\label{tab:confirm2}
\end{table}

\textbf{On a second family and a second benchmark, no baseline is
significantly above PTS.} On \texttt{Qwen2.5-7B-Instruct} the tuning split
selects ablation as the action, and PTS reaches selectivity $0.172$
$[0.136, 0.210]$, $0.081$ above the no-edit arm. After Holm adjustment it is
above the CAST-style trace condition and level with the other six methods,
prompting at $0.212$ included (Table~\ref{tab:confirm2}). On ARC-Challenge PTS
reaches $0.235$ $[0.168, 0.303]$, lowers ECE by $0.060$ and lies above the
95th percentile of a hundred random directions through its decision. It is
above Linear-AcT and the CAST prompt condition, whose decision is at chance on
ARC (AUROC $0.51$), and level with additive CAA and the CAST-style trace
condition, which reach $0.294$ and $0.289$ (Holm $p=0.36$).

\begin{table}[t]\footnotesize\setlength{\tabcolsep}{3pt}
\centering
\resizebox{\linewidth}{!}{%
\begin{tabular}{l|cccc|cccc}
\hline
 & \multicolumn{4}{c|}{Gemma-2-2B ($n_{\mathrm{tox}}{=}451$, $n_{\mathrm{clean}}{=}1{,}910$)} & \multicolumn{4}{c}{Qwen2.5-7B ($n_{\mathrm{tox}}{=}608$, $n_{\mathrm{clean}}{=}1{,}907$)} \\
method & toxic removed (\%) & clean kept (\%) & Sel $\uparrow$ [95\% CI] & $\Delta$Sel ($p$) & toxic removed (\%) & clean kept (\%) & Sel $\uparrow$ [95\% CI] & $\Delta$Sel ($p$) \\
\hline
\inputtab{gen/toxicity}
\hline
\end{tabular}}
\caption{\textbf{PTS decided from the prompt removes most toxic continuations
and keeps most clean ones.} 3{,}000 held-out RealToxicityPrompts prompts per
model, every method at the configuration it won on a separate tuning split. A
clean continuation is kept when the steered one is clean and stays within
sentence-embedding cosine $0.5$ of the unsteered one. $\Delta$Sel is the row
minus PTS with the prompt decision. The last two rows keep that decision and
replace the action by random unit directions at the same dose (mean of 100) or
by a refusal.}
\label{tab:toxicity}
\end{table}

\textbf{On free text, the transport action is above refusing on the same
decision.} Free text has no single readout that an override could rewrite,
so we test PTS on toxic continuation of
RealToxicityPrompts~\citep{gehman2020realtoxicityprompts}, a setting of
Linear-AcT, in two model families. A RoBERTa classifier labels a
continuation toxic at score $0.5$, and a clean continuation is kept only if
the steered one is clean and close to the unsteered one, so rewriting counts
as damage. The action shifts along the difference of mean
activations between toxic and clean continuations, and the decision is a probe
on the prompt, read in one forward pass before generation.

On 3{,}000 held-out prompts per model, PTS reaches selectivity $0.526$ and
$0.534$, and the toxic rate falls from $15.0\%$ to $5.2\%$ on Gemma and from
$20.3\%$ to $4.8\%$ on Qwen (Table~\ref{tab:toxicity}). PTS is above
prompting and the global shift in both families and above CAST on Qwen; on
Gemma it is level with CAST ($+0.011$, $p=0.58$). Refusing on the same
decision converts every clean continuation the decision flags, whereas the
shift converts only part of them, and PTS is above refusing by $0.099$
$[0.066, 0.130]$ and $0.077$ $[0.054, 0.099]$, with the sign that
Eq.~\eqref{eq:override} predicts from tuning-split conversion rates for all six
gated arms. The direction of PTS lies above all hundred random directions
through the same decision on Gemma and above their 95th percentile on Qwen.
Since the decision gates each row inside the original batch, the no-edit arm
leaves all 3{,}000 continuations bitwise unchanged.

\textbf{In the setting of PCHI, a one-direction edit under the same probe is
more selective than a learned head rescaling.} PCHI~\citep{liu2026probe}
gates a learned rescaling of attention heads with a probe read where the model
states its confidence. We rebuild its setting: \texttt{Qwen3-4B-Instruct}
answers OpenMathInstruct-2 problems in JSON ending with a yes/no confidence
field, a probe on the layer-18 state at that field flags overconfident-wrong
answers, and PCHI rescales 128 heads with coefficients learned under its
published hyper-parameters. PTS keeps the probe and replaces the learned rescaling by a
shift of the same state along one direction. On 3{,}183 held-out answers it
reaches selectivity $0.438$ against $0.354$ for PCHI ($+0.084$
$[0.058, 0.110]$), and an expected calibration error of $0.188$ against
$0.234$ (Appendix~\ref{app:pchi}).

\textbf{Deciding from the prompt adds $9\%$ to the unsteered wall clock.}
On one idle H100, 750 toxicity continuations with Gemma take $5.79$ seconds
unsteered and $6.31$ seconds with PTS, of which $0.46$ seconds read the prompt
for the decision; a post-hoc decision reaches $0.723$ and $0.648$
(Table~\ref{tab:toxicity}) but regenerates the flagged half, at $8.74$ seconds.

\section{Limitations}

Our study has several limitations. First, when a behavior is read from a
single output, as multiple-choice confidence is, overriding that output on the
answers the PTS decision flags is at least as selective as the action
($0.319$ against $0.294$ on Gemma and $0.391$ against $0.172$ on Qwen,
Appendix~\ref{app:refs}), and in the setting of PCHI a large random edit at the
readout matches the PTS selectivity in 46 of 100 draws (Appendix~\ref{app:pchi});
the action improves on an override only on free text (Sec.~\ref{sec:transfer}). Second, on Qwen PTS is level with prompting, and its direction is not distinguishable from random
directions displaced as far. The ablation selected there creates about as many
overconfident-wrong answers as it removes, so neither their share nor the
calibration error moves. Third, on Gemma the edit raises the share of
reasoning traces that end without a boxed answer by $5.7$ points; counting only
traces boxed in both arms, the lead over the CAST-style trace condition is $+0.032$
$[-0.008, 0.071]$ (Appendix~\ref{app:net}). Fourth, a single classifier scores
toxicity, and the evidence covers two behaviors in two instruction-tuned
families at the 2B to 7B scale. Fifth, the configuration hashes of
the multiple-choice comparison were committed before its held-out run but made
public only with its results, whereas every later protocol was public before
its runs (Appendix~\ref{app:protocol}). Sixth, the free-text operator reads
the prompt in a pass of its own; editing generated positions only lowers its
selectivity to $0.269$ on Gemma and $0.460$ on Qwen (Appendix~\ref{app:toxicity}).

\section{Conclusion}

We have introduced Projection-Transport Steering, which splits an activation
edit into an action (a transfer function on the projection onto a behavioral
subspace) and a decision (a test on a detection statistic that sets where the
action runs). Once the subspace and the target law are fixed, the cost-minimal action is the optimal-transport map between the
projection laws, and additive steering is that map only when the target is a
translate of the unsteered law. The selectivity of the composed operator equals
the decision's error rates weighted by the action's conversion rates, an
identity that predicts gated configurations from one ungated run; with
Theorem~\ref{thm:gate} it caps, from unsteered answers alone, what a decision
can reach with an action that converts both states alike.

In a comparison committed before its held-out run, PTS is the most selective of
eight methods on MMLU with Gemma, at $0.294$ against $0.248$ for the
strongest baseline, and lowers both the share of overconfident-wrong answers
and the calibration error. On free text, deciding from the prompt before
generating, it is above the global shift and above refusing on the same decision
in two model families.

\subsection*{AI use statement}

We used generative AI tools to polish the manuscript and to retrieve related
work. The authors verified every resulting claim and reference and take
responsibility for the final content.

\providecommand{\bysame}{\leavevmode\hbox to3em{\hrulefill}\thinspace}
\bibliographystyle{iclr2027_conference}
\begin{thebibliography}{10}

\bibitem[Arditi et al.(2024)]{arditi2024refusal}
Andy Arditi, Oscar Obeso, Aaquib Syed, Daniel Paleka, Nina Panickssery, Wes
  Gurnee, and Neel Nanda, \emph{Refusal in language models is mediated by a
  single direction}, Advances in Neural Information Processing Systems 37, 2024.

\bibitem[Azaria \& Mitchell(2023)]{azaria2023lying}
Amos Azaria and Tom Mitchell, \emph{The internal state of an LLM knows when it's
  lying}, Findings of the Association for Computational Linguistics: EMNLP, 2023.

\bibitem[Braun et al.(2025)]{braun2025unreliability}
J.~Braun, C.~Eickhoff, D.~Krueger, S.~A.~Bahrainian, D.~Krasheninnikov.
  \emph{Understanding (un)reliability of steering vectors in language models}.
  ICLR Workshop on Foundation Models in the Wild, 2025.

\bibitem[Clark et al.(2018)]{clark2018arc}
Peter Clark, Isaac Cowhey, Oren Etzioni, Tushar Khot, Ashish Sabharwal, Carissa
  Schoenick, and Oyvind Tafjord, \emph{Think you have solved question answering?
  Try ARC, the AI2 reasoning challenge}, arXiv:1803.05457, 2018.

\bibitem[Cobbe et al.(2021)]{cobbe2021gsm8k}
Karl Cobbe, Vineet Kosaraju, Mohammad Bavarian, Mark Chen, Heewoo Jun, Lukasz
  Kaiser, et al., \emph{Training verifiers to solve math word problems},
  arXiv:2110.14168, 2021.

\bibitem[Fan et al.(2026)]{steerable2026}
Chenrui Fan, Yize Cheng, Ming Li, Soheil Feizi, and Tianyi Zhou, \emph{When is
  your LLM steerable? Predicting steering success from internal states},
  arXiv:2606.11599, 2026.

\bibitem[Ferrando et al.(2026)]{ferrando2026dsas}
Javier Ferrando, Xavier Suau, Pau Rodriguez, and Adri\`a Gonzalez,
  \emph{Dynamically scaled activation steering}, arXiv:2512.03661, 2026.

\bibitem[Gemma Team(2024)]{gemma2024}
Gemma Team, \emph{Gemma 2: Improving open language models at a practical size},
  arXiv:2408.00118, 2024.

\bibitem[Haider et al.(2026)]{gaps2026}
Umair Haider, Hassan Sajjad, and A.~B. Siddique, \emph{GAPS: gated activation
  steering at dimension granularity}, arXiv:2609.01878, 2026.

\bibitem[Gehman et al.(2020)]{gehman2020realtoxicityprompts}
Samuel Gehman, Suchin Gururangan, Maarten Sap, Yejin Choi, and Noah~A. Smith,
  \emph{RealToxicityPrompts: Evaluating neural toxic degeneration in language
  models}, Findings of the Association for Computational Linguistics: EMNLP, 2020.

\bibitem[Gelbrich(1990)]{gelbrich1990formula}
Matthias Gelbrich, \emph{On a formula for the $L^2$ Wasserstein metric between
  measures on Euclidean and Hilbert spaces}, Mathematische Nachrichten
  \textbf{147} (1990), no.~1, 185--203.

\bibitem[Hendrycks et al.(2021)]{hendrycks2021mmlu}
Dan Hendrycks, Collin Burns, Steven Basart, Andy Zou, Mantas Mazeika, Dawn Song,
  and Jacob Steinhardt, \emph{Measuring massive multitask language
  understanding}, International Conference on Learning Representations, 2021.

\bibitem[Kadavath et al.(2022)]{kadavath2022know}
Saurav Kadavath, Tom Conerly, Amanda Askell, Tom Henighan, Dawn Drain, Ethan
  Perez, et al., \emph{Language models (mostly) know what they know},
  arXiv:2207.05221, 2022.

\bibitem[Knott \& Smith(1984)]{knott1984optimal}
Martin Knott and Cyril~S. Smith, \emph{On the optimal mapping of distributions},
  Journal of Optimization Theory and Applications \textbf{43} (1984), 39--49.

\bibitem[Kortukov et al.(2026)]{fpcg2026}
Evgenii Kortukov, Piotr Komorowski, Florian Klein, Paula Engl, Gabriele Sarti,
  Seong Joon Oh, Sebastian Lapuschkin, and Wojciech Samek, \emph{Predicting
  future behaviors in reasoning models enables better steering},
  arXiv:2606.11172, 2026.

\bibitem[Lee et al.(2025)]{lee2024cast}
Bruce~W. Lee, Inkit Padhi, Karthikeyan~Natesan Ramamurthy, Erik Miehling,
  Pierre Dognin, Manish Nagireddy, and Amit Dhurandhar, \emph{Programming
  refusal with conditional activation steering}, International Conference on
  Learning Representations, 2025.

\bibitem[Li et al.(2023)]{li2023iti}
Kenneth Li, Oam Patel, Fernanda Vi\'egas, Hanspeter Pfister, and Martin
  Wattenberg, \emph{Inference-time intervention: Eliciting truthful answers from
  a language model}, Advances in Neural Information Processing Systems 36, 2023.

\bibitem[Li et al.(2026)]{liu2026probe}
Ke Li, Chongzhe Zhang, Zifan Zeng, Feng Liu, Qunli Zhang, and Zheng Hu,
  \emph{Calibrating overconfidence without sacrificing confidence:
  Probe-conditioned head intervention for LLMs}, arXiv:2606.09876, 2026.

\bibitem[Luo et al.(2026)]{luo2026yopo}
Ziyang Luo, Zhongyao Chu, Xinjie He, Youting Wang, Xukui Qin, Runxiong Wu, and
  Yan-Syuan Chen, \emph{You only pass once: answering and abstaining together
  in a single forward pass of a frozen language model}, arXiv:2608.14465, 2026.

\bibitem[Nguyen et al.(2026)]{nguyen2026coast}
Tam Nguyen, Tu~Anh Nguyen, Sina Alemohammad, and Richard~G. Baraniuk,
  \emph{Minimizing collateral damage in activation steering}, arXiv:2605.01167,
  2026.

\bibitem[Panickssery et al.(2024)]{rimsky2024caa}
Nina Panickssery, Nick Gabrieli, Julian Schulz, Meg Tong, Evan Hubinger, and
  Alexander~Matt Turner, \emph{Steering Llama 2 via contrastive activation
  addition}, Proceedings of the 62nd Annual Meeting of the Association for
  Computational Linguistics, 2024.

\bibitem[Rodriguez et al.(2025a)]{rodriguez2025act}
Pau Rodriguez, Arno Blaas, Michal Klein, Luca Zappella, Nicholas Apostoloff,
  Marco Cuturi, and Xavier Suau, \emph{Controlling language and diffusion models
  by transporting activations}, International Conference on Learning
  Representations, 2025.

\bibitem[Rodriguez et al.(2025b)]{rodriguez2025lineas}
Pau Rodriguez, Michal Klein, Eleonora Gualdoni, Arno Blaas, Luca Zappella, Marco
  Cuturi, and Xavier Suau, \emph{LinEAS: End-to-end learning of activation
  steering with a distributional loss}, arXiv:2503.10679, 2025.

\bibitem[Scalena et al.(2024)]{scalena2024dynamic}
Daniel Scalena, Gabriele Sarti, and Malvina Nissim, \emph{Multi-property
  steering of language models with dynamic activation composition}, Proceedings
  of BlackboxNLP, 2024.

\bibitem[Singh et al.(2024)]{singh2024mimic}
Shashwat Singh, Shauli Ravfogel, Jonathan Herzig, Roee Aharoni, Ryan Cotterell,
  and Ponnurangam Kumaraguru, \emph{Representation surgery: Theory and practice
  of affine steering}, Proceedings of the 41st International Conference on
  Machine Learning, 2024.

\bibitem[Skifstad et al.(2026)]{nguyen2026alqr}
Julian Skifstad, Xinyue~Annie Yang, and Glen Chou, \emph{Local Linearity of
  LLMs Enables Activation Steering via Model-Based Linear Optimal Control},
  arXiv:2604.19018, 2026.

\bibitem[Tan et al.(2024)]{tan2024generalisation}
D.~Tan, D.~Chanin, A.~Lynch, B.~Paige, D.~Kanoulas, A.~Garriga-Alonso, R.~Kirk.
  \emph{Analysing the generalisation and reliability of steering vectors}.
  NeurIPS, 2024.

\bibitem[Templeton et al.(2024)]{templeton2024scaling}
Adly Templeton, Tom Conerly, Jonathan Marcus, et al., \emph{Scaling
  monosemanticity: Extracting interpretable features from Claude 3 Sonnet},
  Transformer Circuits Thread, 2024.

\bibitem[Turner et al.(2023)]{turner2023actadd}
Alexander~Matt Turner, Lisa Thiergart, Gavin Leech, David Udell, Juan~J.
  Vazquez, Ulisse Mini, and Monte MacDiarmid, \emph{Steering language models
  with activation engineering}, arXiv:2308.10248, 2023.

\bibitem[Wang et al.(2025)]{wang2025sadi}
Weixuan Wang, Jingyuan Yang, and Wei Peng, \emph{Semantics-adaptive activation
  intervention for LLMs via dynamic steering vectors}, International Conference
  on Learning Representations, 2025.

\bibitem[Wu et al.(2025)]{wu2025axbench}
Zhengxuan Wu, Aryaman Arora, Atticus Geiger, Zheng Wang, Jing Huang, Dan
  Jurafsky, Christopher~D. Manning, and Christopher Potts, \emph{AxBench:
  Steering LLMs? Even simple baselines outperform sparse autoencoders},
  Proceedings of the 42nd International Conference on Machine Learning, 2025.

\bibitem[Zou et al.(2023)]{zou2023repe}
Andy Zou, Long Phan, Sarah Chen, James Campbell, et al., \emph{Representation
  engineering: A top-down approach to AI transparency}, arXiv:2310.01405, 2023.

\end{thebibliography}

\appendix

\section{Confirmatory protocol}
\label{app:protocol}

The multiple-choice protocol was committed to the local repository before any
steered rollout on the confirmatory split was generated and first pushed
together with the results, so its ordering rests on local commit times; every
later protocol below was pushed to a public repository, and so timestamped by
the hosting service, before its rollouts existed (\texttt{rebuttal/protocols/confirmatory-v19.md},
with two amendments made after the tuning pass and before launch), together
with the SHA-256 hashes of the selected configurations and of every fitted
detector (\texttt{confirmatory-v19-gemma.sha256}).

\textbf{Splits.} All splits are subsets of the 14{,}042-question MMLU test set,
identified by index and constructed by \texttt{label\_pool.split\_records}. The
extraction split (400) excludes the first evaluation subset of the pilot
protocol (Appendix~\ref{app:legacy}); the tuning split (1{,}200) excludes the
extraction split and the five pilot evaluation subsets; the detector pool (4{,}000) and the confirmatory
split (3{,}000) are consecutive blocks of a seeded permutation of the
remaining questions. No question appears in two of the four splits.

\textbf{Selection rule.} For every method the configuration with the highest
tuning-split selectivity among those with $\Delta\mathrm{acc}\ge-0.01$ is
carried unchanged to the confirmatory split. Gated configurations are scored on
the tuning split by combining the ungated rollouts of the same split with the
gate's flags (Prop.~\ref{prop:account} makes this exact up to batched-decoding
nondeterminism; Appendix~\ref{app:law}); single-pass configurations are real
runs. Every confirmatory number is a real run.

\textbf{Tests.} Primary: PTS post-hoc against each of the seven baselines on
pooled M5 selectivity, question-level paired bootstrap with $10^4$ resamples,
two-sided, Holm-adjusted over the seven comparisons. Accuracy: an arm is
non-inferior if the lower end of the $95\%$ interval of its accuracy change is
above $-0.02$. The single-pass and prompt-decision variants, M4 and M2
selectivities, removal, retention and ECE are secondary and reported without
a claim.

\textbf{Released records.} Every unsteered and steered answer of the tuning and
confirmatory splits is released with its state, confidence and correctness
(\texttt{results/release/records.csv.gz}); \texttt{paper/recompute\_from\_records.py}
recomputes each selectivity and paired interval in Tables~\ref{tab:confirm}
and~\ref{tab:confirm2}, including the null arms, from that file alone.
\texttt{results/release/toxicity.csv.gz} holds every confirmatory continuation
of the free-text setting with its judge score and state (texts omitted for the
random-direction arms), and \texttt{results/release/pchi\_val.csv.gz} the
held-out readout gaps of every arm of the PCHI setting.

\textbf{ARC replication.} A second protocol
(\texttt{confirmatory-v20-arc.md}), frozen before any steered ARC rollout,
repeats the Gemma comparison on ARC-Challenge: a detector pool of 1{,}012 and a
tuning split of 400 from the ARC training and validation questions, and the
1{,}165-question test split as the confirmatory split. The directions, moments
and action layer are those fitted on MMLU; every decision is refitted on the ARC
pool together with the MMLU pool, and every configuration is selected on the
ARC tuning split with the same rule and grids.

\textbf{Null arms.} A third protocol (\texttt{confirmatory-v21-null.md}), with
the job list and the hash of the runner, was committed and pushed to a public
repository before any null rollout existed; the commit is timestamped by the
hosting service. On each confirmatory split it adds the frozen PTS decision with
an identity action, the frozen decision with the selected action on a unit
random direction (seed 0), and an ungated rerun with an identity action. It
registers two one-sided hypotheses per setting, PTS above each gated null, with
Holm adjustment over the six tests, and the net measures of
Appendix~\ref{app:net} for every arm. Five of the six tests pass; PTS against
the single random direction on ARC gives $+0.041$ $[-0.029, 0.109]$, Holm
$p=0.12$.

\textbf{Reference arms.} A fourth protocol
(\texttt{confirmatory-v22-references.md}), pushed before any of its rollouts
existed, extends the no-edit null to all four conditional methods, replaces
the single random direction by ten (seeds 0 to 9), and adds random directions
for the sycophancy edit. It registers 24 one-sided tests (four methods, three
settings, no-edit and random-edit null), Holm-adjusted together, of which 18
pass (Table~\ref{tab:refs}); a count of the cases in which
Eq.~\eqref{eq:override} with ungated conversion rates predicts the sign of
steer minus matched override (10 of 13); and paired intervals for selectivity
at cutoffs 0.4 and 0.6 and over boxed traces for PTS against the CAST-style
condition on Gemma. The v22 random-edit null supersedes the single-direction
test of v21.

\textbf{Free text and dense nulls.} A fifth protocol
(\texttt{confirmatory-v23-freetext.md}), pushed with the hashes of the tuning
selections, fitted decisions and code before any confirmatory rollout of the
toxicity setting existed, fixes the data, judge, retention criterion, arms and
ten primary hypotheses of Appendix~\ref{app:toxicity}, and extends the random
directions of PTS post-hoc to one hundred on Gemma-MMLU and Gemma-ARC and to
thirty displacement-matched directions on Qwen (Table~\ref{tab:dense}).

\textbf{PCHI setting.} A sixth protocol (\texttt{confirmatory-v24-pchi.md}),
pushed after the train-split fitting and tuning and before any validation
rollout, fixes the arms and three hypotheses of Appendix~\ref{app:pchi}.

\textbf{Analyses added after the protocol.} The detect-and-override reference,
the difference-of-means decision fitted on 4{,}400 labels, the latency
measurement (Appendix~\ref{app:latency}), the decision-by-action
matrix and the trace-only and confidence-turn-only edits
(Appendix~\ref{app:decisions}), the sycophancy experiment
(Appendix~\ref{app:sycophancy}), calibration and the out-of-domain negatives were run
after the confirmatory results were known. They use the same splits and the
frozen decisions, and we report them as exploratory.

\section{A pilot operator on further models}
\label{app:postreview}

In a pilot study we froze a simplified operator, an action along a
detection direction, after a development analysis on Qwen3, and tested it on
held-out TruthfulQA panels of 237 questions (50 in-domain, 187 shifted) for
Mistral-7B, OLMo-2-7B and Granite-8B. Each comparison was against a
norm-matched spherical action, with 10{,}000 paired BCa resamples and
Bonferroni-adjusted two-sided intervals, and a model passed only if calibrated
MC2 and MC1 both improved with intervals excluding zero.

\begin{table}[h]\footnotesize\setlength{\tabcolsep}{3pt}
\centering
\resizebox{\linewidth}{!}{%
\begin{tabular}{llrrl}
\hline
model / panel & comparison & $\Delta$ calibrated MC2 [adj. 95\% CI] & $\Delta$ MC1 [adj. 95\% CI] & decision \\
\hline
OLMo-2 / TruthfulQA & matched spherical & $+.137\ [+.105,+.168]$ & $+.105\ [+.025,+.177]$ & pass \\
OLMo-2 / TruthfulQA & no intervention & $+.299\ [+.251,+.345]$ & $+.287\ [+.177,+.384]$ & pass \\
OLMo-2 / TruthfulQA & strongest isotropic & $+.220\ [+.180,+.258]$ & $+.194\ [+.097,+.287]$ & pass \\
OLMo-2 / TruthfulQA & unprompted fit & $+.295\ [+.247,+.343]$ & $+.287\ [+.177,+.384]$ & pass \\
\hline
\end{tabular}}
\caption{Held-out OLMo-2 validation on 237 TruthfulQA questions, 50 in-domain
and 187 shifted. Intervals use 10,000 paired BCa resamples and the predeclared
Bonferroni family; the table shows the matched action and representative
construction controls of the seven-control family, all of which pass.}
\label{tab:postreview}
\end{table}

On OLMo-2 the operator passed against the matched action and against all seven
pre-registered construction controls (Table~\ref{tab:postreview}). Mistral
improved calibrated MC2 by $+0.153$ $[+0.108,+0.197]$, but its MC1 interval
$[-0.059,+0.122]$ included zero, and Granite did not clear the screening
threshold. A separate screen of the same three models on BBQ and ETHICS, six
model--behavior cells with disjoint fitting, screening and test units, advanced
one cell, Mistral on ETHICS, to a 1{,}000-scenario test. There the operator
improved pooled calibrated MC2 by $+0.107$ $[+0.090,+0.124]$ and lowered MC1 on
label-1 scenarios by $0.093$ $[0.071,0.124]$.

\section{Proofs}
\label{app:proofs}

\textbf{Theorem~\ref{thm:ot}.} By locality, $T(h)-h\in\mathrm{span}(V^\top)$
$\mu$-a.s., hence $\lVert T(h)-h\rVert^2=\lVert VT(h)-Vh\rVert^2$ because the
rows of $V$ are orthonormal, and $C(T)=\mathbb{E}\lVert VT(h)-Vh\rVert^2$. The
joint law $\gamma$ of $(Vh,VT(h))$ is a coupling of $(\mu_V,\nu)$ by the goal
constraint, so $C(T)=\int\lVert y-x\rVert^2d\gamma\ge\Wtwo^2(\mu_V,\nu)$.
Because $\mu_V$ is absolutely continuous and both laws have finite second
moments, Brenier's theorem gives a unique optimal coupling, induced by
$S=\nabla\varphi$ with $\varphi$ convex. $T^*(h)=h+V^\top(S(Vh)-Vh)$ is
admissible ($VV^\top=I_k$ gives $VT^*(h)=S(Vh)$, and the complement is
untouched) and attains the bound. For any minimiser the induced coupling is
optimal, hence equal to the unique one, so $VT(h)=S(Vh)$ a.s.\ \qed

For $k=1$ absolute continuity can be weakened to atomlessness of $\mu_V$: the
monotone rearrangement $F_\nu^{-1}\circ F_{\mu_V}$ is then an optimal map for
every $\nu$ with a finite second moment, including targets with atoms such as
$\delta_0$ (ablation) or a law truncated at $c$ (clamp); uniqueness of the map
can fail only on a $\mu_V$-null set.

\begin{proposition}[metric covariance]
\label{prop:metric}
If the cost is $(\Delta h)^\top M(\Delta h)$ with $M\succ0$, the admissible
problem is optimal transport of $\mu_V$ to $\nu$ in $\mathbb{R}^k$ under the
cost $\lVert\cdot\rVert^2_{M_V}$, $M_V=VMV^\top\succ0$, and the optimal
operator is $T^*(h)=h+V^\top(S_M(Vh)-Vh)$ with
$S_M=M_V^{-1/2}\circ\nabla\varphi\circ M_V^{1/2}$, $\varphi$ convex.
\end{proposition}

\emph{Proof.} Locality writes $\Delta h=V^\top a$ with $a\in\mathbb{R}^k$, and
the goal constraint with $VV^\top=I_k$ fixes $a=S(Vh)-Vh$ for the map $S$ in
use. Then $\lVert\Delta h\rVert_M^2=a^\top M_Va$. The cost is a latent-space
Mahalanobis cost on $\mathbb{R}^k$; it is not the full-space Mahalanobis
distance of $\Delta h$ unless $M$ commutes with $V^\top V$. In the coordinates
$z=M_V^{1/2}q$ it is Euclidean, the change of variables is internal to
$\mathbb{R}^k$, and Brenier's theorem in $z$ (the pushforward of an absolutely
continuous law under an invertible linear map is absolutely continuous) gives
$\nabla\varphi$; mapping back gives $S_M$. \qed

\textbf{Corollary~\ref{cor:additive}.} (a) If $\nu=\mu_V*\delta_\alpha$ the
map $x\mapsto x+\alpha$ is the gradient of a convex function pushing $\mu_V$
to $\nu$, hence the Brenier map. Conversely, if the optimal map is
$x+\alpha$ then $\nu$ is its pushforward, the shift. For one-dimensional
Gaussians the monotone map is $x\mapsto m_t+(\sigma_t/\sigma_s)(x-m_s)$, with
slope $\sigma_t/\sigma_s$. (b) If two groups of inputs have the same law of $p$
and the probability that an input changes state depends on $h$ only through
$p$, every transfer function changes both groups at the same rate, since both
rates are the same integral against the law of $p$. A decision that reads $p$
alone therefore adds little, which is why the detection statistic reads the
full activation. \qed

\textbf{Proposition~\ref{prop:account}.} Let $R$ be the event that an \ocw{}
trace leaves \ocw{} and $L$ the event that a \crok{} trace leaves \crok{}. Off
$\Gamma$ the output equals the unsteered output, which is in its original state
by definition, so $R\subseteq\Gamma$ and $L\subseteq\Gamma$ up to null sets.
Then $P(R\mid\ocw)=P(\Gamma\mid\ocw)P(R\mid\ocw,\Gamma)=\mathrm{TPR}\,\rho_O$, and
likewise $P(L\mid\crok)=\mathrm{FPR}\,\rho_C$; subtracting gives
Eq.~\eqref{eq:account}. If, within \ocw{}, the indicator that the action
converts the trace is independent of $\Gamma$, then
$\rho_O=P(R^{\mathrm{ungated}}\mid\ocw)$, and likewise for \crok{}; the ungated
selectivity is $\rho_O-\rho_C$, and the difference is
$\rho_C(1-\mathrm{FPR})-\rho_O(1-\mathrm{TPR})$. \qed

The identity requires that the unflagged output be exactly the unsteered one.
The post-hoc gate satisfies this by construction. The single-pass gate
satisfies it for traces on which it never fires, since the action hook is the
identity until the decision token.

\textbf{Theorem~\ref{thm:gate}.} Write $p_O,p_C$ for the densities of $s$ on
\ocw{} and \crok{} answers and $r_O(s),r_C(s)$ for the probabilities that the
action converts an answer of each state with statistic $s$. For a decision
region $\Gamma$, $\mathrm{TPR}\,\rho_O=\int_\Gamma p_O r_O$ and
$\mathrm{FPR}\,\rho_C=\int_\Gamma p_C r_C$, so
$\mathrm{Sel}(\Gamma)=\int_\Gamma(p_O r_O-p_C r_C)$. Among regions with a given
value of $\int_\Gamma p_C r_C$, the Neyman--Pearson lemma applied to the
nonnegative measures $p_O r_O$ and $p_C r_C$ shows that
$\int_\Gamma p_O r_O$ is maximised by a superlevel set of their ratio, which is
the general statement in (a). If $r_O$ and $r_C$ are constant in $s$, the ratio
is a constant multiple of $p_O/p_C$, so every objective increasing in TPR and
decreasing in FPR is maximised on the likelihood-ratio ROC curve, which is (a).
The assumption matters: with $s\in\{a,b\}$, $p_O(a)=p_O(b)=\tfrac12$,
$p_C(a)=\tfrac14$, $p_C(b)=\tfrac34$, $r_O(a)=0.2$, $r_O(b)=1$, $r_C(a)=1$ and
$r_C(b)=0$, the likelihood ratio ranks $a$ above $b$, yet $\Gamma=\{b\}$ gives
selectivity $0.5$, against $-0.15$ for $\{a\}$ and $0.35$ for $\{a,b\}$. On
our data $r$ does depend on $s$: flagged answers are converted more often than
unflagged answers of the same state (Appendix~\ref{app:law}). Our detector is a
logistic-regression probe, a discriminative estimate of the log-likelihood
ratio up to the prior term; (a) states what the probe approximates, not that
the fitted probe is optimal. (b) The concavified ROC contains the path
$(0,0)\to(\mathrm{FPR}^*,\mathrm{TPR}^*)\to(1,1)$ through the Youden-optimal
point, whose area is $\tfrac12(1+J)$, so $J\le2\,\mathrm{AUC}-1$. The override
reference has selectivity $\mathrm{TPR}-\mathrm{FPR}\le J$, and with
$\rho_O=\rho_C=\rho$, Eq.~\eqref{eq:account} gives $\mathrm{Sel}=\rho(\mathrm{TPR}-\mathrm{FPR})\le J$. \qed

\textbf{Theorem~\ref{thm:joint}.} Split $C(T)$ over $\Gamma$ and $\Gamma^c$.
The second term is zero by the identity constraint; the first, with the
conditional goal constraint, is Theorem~\ref{thm:ot} for
$\mu(\cdot\mid\Gamma)$, which applies when $\mu_V(\cdot\mid\Gamma)$ is
absolutely continuous. \qed

This result fixes the decision. It does not say that a threshold gate is the
best decision for the displacement-priced objective; the next statement shows
that it generally is not, and when it is.

\begin{proposition}[graded dose]
\label{prop:dose}
Let $T_\lambda(h)=h+\lambda(s)\,V^\top(S(Vh)-Vh)$ with a dose
$\lambda:\mathbb{R}\to[0,1]$ measurable in the detector statistic $s$. Write
$\pi(s)=P(\ocw\mid s)$, $d(s)=\mathbb{E}[\lVert S(Vh)-Vh\rVert^2\mid s]>0$, and
suppose the behavioral gain is linear in the dose, with slope $\beta_1>0$ on
\ocw{} and $-\beta_0<0$ on \crok{}. The dose maximising gain minus
displacement priced at $1/(2c)$ is
\[
\lambda^*(s)=\Big[\tfrac{c\,(\beta_0+\beta_1)\,(\pi(s)-\tau)}{d(s)}\Big]_0^1,\qquad
\tau=\tfrac{\beta_0}{\beta_0+\beta_1}.
\]
It is zero exactly on $\{\pi\le\tau\}$ and nondecreasing in the ratio
$(\pi(s)-\tau)/d(s)$. It is nondecreasing in $\pi$ when $d$ is constant, but
not in general. The threshold gate $\mathbf{1}[\pi>\tau]$ is optimal iff
$\lambda^*\in\{0,1\}$ almost surely, and strictly suboptimal iff
$P(0<\lambda^*<1)>0$.
\end{proposition}

\emph{Proof.} The displacement of $T_\lambda$ is
$\lambda(s)^2\lVert S(Vh)-Vh\rVert^2$, with expectation
$\mathbb{E}_s[\lambda(s)^2d(s)]$; the gain is
$\mathbb{E}_s[\lambda(s)((\beta_0+\beta_1)\pi(s)-\beta_0)]$. The objective is
an integral over $s$ of a strictly concave quadratic in $\lambda$, maximised
pointwise on $[0,1]$ by the clipped stationary point. It is the unique
maximiser a.s., so the indicator is optimal iff it agrees with $\lambda^*$
a.s., i.e.\ iff $\lambda^*\in\{0,1\}$ a.s.; otherwise the objective is strictly
larger at $\lambda^*$ on a set of positive measure. For non-monotonicity in
$\pi$, take $s_1,s_2$ with $\pi(s_1)=0.6<\pi(s_2)=0.7$, $\tau=0.5$,
$c(\beta_0+\beta_1)=1$, $d(s_1)=0.2$ and $d(s_2)=1$: then
$\lambda^*(s_1)=0.5>\lambda^*(s_2)=0.2$. For a counterexample to strict
suboptimality at every finite price, let $\pi$ take only the values $0.1$ and
$0.9$ with $\tau=0.5$ and $d\le0.4\,c(\beta_0+\beta_1)$: then
$\lambda^*\in\{0,1\}$ a.s.\ and the threshold gate is optimal. \qed

\textbf{Proposition~\ref{prop:seq}.} Let $\mathcal{F}_t$ be generated by the first $t$
generated tokens and $z_t$ the mean detector feature over them. A decision
measurable with respect to $\sigma(z_t)$ is a test between the
class-conditional laws of $z_t$, so Theorem~\ref{thm:gate} applies with $z_t$
in place of $s$ and nothing is assumed about how $z_t$ relates to the finished
trace. The reader hook sits at a layer at or above the actor; within one
forward pass the actor at position $t$ sees the decision accumulated over
positions $<t$, the running mean needs one sum and one counter per sequence,
and an edit at position $t$ changes only the keys and values written at $t$.

\begin{proposition}[computational cost]
\label{prop:compute}
Any algorithm implementing a non-degenerate intervention local to a subspace
spanned by a dense direction reads all $d$ coordinates of $h$, so it costs
$\Omega(d)$ per token. The edit of Eq.~\eqref{eq:edit} costs $\Theta(kd+k\log K)$ with a
$K$-knot lookup, and the detector adds one $d$-dimensional dot product per
token.
\end{proposition}

\emph{Proof.} If a coordinate $i$ with $(\hat v)_i\neq0$ is never read, an
adversary perturbs $h_i$, which changes $p=\langle h,\hat v\rangle$ and, on a
set of positive measure, the correct output, while the algorithm's output is
unchanged. The upper bound counts $Vh$, a binary search in the lookup, and
$h+V^\top(\cdot)$. \qed

\section{GSM8K MCQ construction}
\label{app:gsm}
Each GSM8K test item with integer gold $g$ receives three deterministic
distractors drawn (seeded by item index) from
$\{g\pm1, g\pm2, g\pm10, 2g, \lfloor g/2\rfloor, \mathrm{round}(1.5g),
\mathrm{round}(0.8g), g\pm\mathrm{U}\{3..9\}\}$, deduplicated, shuffled with the
gold into four options. This converts free-form arithmetic into selection while
keeping the reasoning burden; we report transfer numbers on it and do not claim
GSM8K free-form results.

\section{Tuning grids}
\label{app:tuning}

Every configuration searched on the tuning split, with the selected one marked
$\star$. Gated rows are scored by composing the ungated rollouts of the same
split with the decision's flags (Appendix~\ref{app:law}); single-pass rows are
real runs. The rule is the highest selectivity among configurations with
$\Delta\mathrm{acc}\ge-0.01$, and, when no configuration of a family meets the
floor, the highest selectivity overall; this fallback applies to the additive,
CAST and prompt-decision families on Qwen, where every shift in the grid costs
more than one accuracy point.

{\scriptsize\setlength{\tabcolsep}{3pt}
\begin{longtable}{l|l|c|cc|c}
\caption{Tuning split, \texttt{gemma-2-2b-it} ($n=1{,}200$, M5).}\label{tab:tune-gemma}\\
\hline
family & configuration & $\Delta$acc & removal & retention & Sel \\
\hline
\endfirsthead
\hline
family & configuration & $\Delta$acc & removal & retention & Sel \\
\hline
\endhead
\inputtab{gen/tuning_gemma}
\hline
\end{longtable}}

{\scriptsize\setlength{\tabcolsep}{3pt}
\begin{longtable}{l|l|c|cc|c}
\caption{Tuning split, \texttt{Qwen2.5-7B-Instruct} ($n=1{,}200$, M5).}\label{tab:tune-qwen}\\
\hline
family & configuration & $\Delta$acc & removal & retention & Sel \\
\hline
\endfirsthead
\hline
family & configuration & $\Delta$acc & removal & retention & Sel \\
\hline
\endhead
\inputtab{gen/tuning_qwen}
\hline
\end{longtable}}

\section{The selectivity identity in practice}
\label{app:law}

Proposition~\ref{prop:account} needs the output to be unchanged wherever the
decision does not fire. For the post-hoc and prompt decisions this holds by
construction, so a gated arm equals, question by question, either its unsteered
row or its ungated row. Composing the two requires one ungated run per action
and one set of decision scores, and it differs from a real gated run only
through batched-decoding nondeterminism: a flagged question is generated in a
different batch in the two cases. Table~\ref{tab:law-confirm} compares the
composition with the real run for the four gated arms of each confirmatory
split, and reports the prediction from the ungated conversion rates
(independence within each state).

\begin{table}[h]\footnotesize\setlength{\tabcolsep}{4pt}
\centering
\begin{tabular}{ll|cc|cc|ccc}
\hline
setting & arm & TPR & FPR & $\rho_O$ & $\rho_C$ & predicted & composed & real run \\
\hline
\inputtab{gen/law_confirm}
\end{tabular}
\caption{Real gated runs against the composition of ungated runs and against
the prediction $\mathrm{TPR}\,\rho_O-\mathrm{FPR}\,\rho_C$ with ungated
conversion rates.}
\label{tab:law-confirm}
\end{table}

Across the twelve arms the composition differs from the real run by at most
$0.025$ (Qwen, PTS post-hoc) and by $0.005$ at the median, so composed cells in
Fig.~\ref{fig:decides} carry an error of that size. The prediction is higher
than the real run in eleven of twelve rows, and the same bias appears on the
tuning split (Fig.~\ref{fig:law}a).
Flagged answers are converted more often than the average answer of their state
in both states: for PTS post-hoc, $0.91$ of flagged overconfident-wrong answers
change state against $0.84$ ungated, and $0.89$ of flagged confident-right
answers against $0.59$ ungated. The larger increase on the protected state
lowers selectivity below the independence prediction. The composition, which
uses no such assumption, is what scores the gated configurations during
selection.

\section{Secondary readouts}
\label{app:secondary}

Table~\ref{tab:secondary} reports the Gemma confirmatory arms under the two
secondary readouts. The letter-logit confidence M2 assigns confidence above
$0.5$ to nearly every answer, so its \ocw{} stratum is close to the set of wrong
answers and every method moves it little. Under M4 the edits act on four
separate verification generations; PTS is level with the ungated edits there,
which is expected for a decision fitted to M5 labels.

\begin{table}[h]\footnotesize\setlength{\tabcolsep}{3pt}
\centering
\resizebox{\linewidth}{!}{%
\begin{tabular}{l|cc|cc}
\hline
 & \multicolumn{2}{c|}{M2 (letter logit)} & \multicolumn{2}{c}{M4 (verification)} \\
method (Gemma, confirmatory) & Sel [95\% CI] & $\Delta$ vs PTS & Sel [95\% CI] & $\Delta$ vs PTS \\
\hline
\inputtab{gen/secondary_gemma}
\end{tabular}}
\caption{Secondary readouts on the confirmatory split; paired bootstrap
intervals, no multiplicity adjustment.}
\label{tab:secondary}
\end{table}

\begin{table}[h]\footnotesize\setlength{\tabcolsep}{3pt}
\centering
\resizebox{\linewidth}{!}{%
\begin{tabular}{l|cc|cc}
\hline
 & \multicolumn{2}{c|}{M2 (letter logit)} & \multicolumn{2}{c}{M4 (verification)} \\
method (Qwen, confirmatory) & Sel [95\% CI] & $\Delta$ vs PTS & Sel [95\% CI] & $\Delta$ vs PTS \\
\hline
\inputtab{gen/secondary_qwen}
\end{tabular}}
\caption{Secondary readouts on the Qwen confirmatory split.}
\label{tab:secondary-qwen}
\end{table}

Table~\ref{tab:transitions} gives the full state transitions of the frozen
PTS arms under M5. Selectivity counts only the first and last rows; the table
also shows the answers the edit corrects or breaks and the few answers it turns
overconfident.

\begin{table}[h]\footnotesize\setlength{\tabcolsep}{5pt}
\centering
\begin{tabular}{ll|cccc}
\hline
 & before $\backslash$ after & OCW & NCW & NCR & CR \\
\hline
\inputtab{gen/transitions}
\end{tabular}
\caption{State transitions of the frozen PTS arms on the confirmatory split
(counts of questions).}
\label{tab:transitions}
\end{table}

Table~\ref{tab:calibration} reports calibration of the stated confidence on
the confirmatory split: expected calibration error (10 bins), Brier score, and
the AUROC with which the confidence separates right from wrong answers.

\begin{table}[h]\footnotesize\setlength{\tabcolsep}{4pt}
\centering
\begin{tabular}{l|ccc|ccc}
\hline
 & \multicolumn{3}{c|}{Gemma-2-2B} & \multicolumn{3}{c}{Qwen2.5-7B} \\
arm & ECE $\downarrow$ & Brier $\downarrow$ & AUROC $\uparrow$ & ECE $\downarrow$ & Brier $\downarrow$ & AUROC $\uparrow$ \\
\hline
\inputtab{gen/calibration}
\hline
\end{tabular}
\caption{Calibration of the stated confidence (M5) on the confirmatory split.}
\label{tab:calibration}
\end{table}

\section{Net effects of every arm}
\label{app:net}

Selectivity is computed on the answers that start in \ocw{} or \crok{}; it does
not see answers that enter \ocw{} from another state. Table~\ref{tab:net}
reports, for every confirmatory arm and the three null arms of
Appendix~\ref{app:protocol}, selectivity net of the no-edit null, the change in
the share of \ocw{} answers among all questions, the change in the rate of
confident answers among wrong and among right answers, $\Delta$ECE,
$\Delta$Brier, the change in the share of traces without a \verb|\boxed{}|
answer (which then receive a forced box, Appendix~\ref{app:prompts}), and
selectivity recomputed without the questions whose steered trace has no box.
The ungated rerun reproduces the unsteered outputs exactly, because it uses the
same batches; the gated no-edit null regenerates the flagged questions in
different batches, and the resulting nondeterminism alone yields selectivity
$0.049$ (Gemma), $0.091$ (Qwen) and $0.016$ (ARC). Excluding traces without a
box leaves the ordering on Gemma unchanged: PTS post-hoc $0.257$, CAST-style
$0.233$, additive $0.215$.

\begin{table}[h]\scriptsize\setlength{\tabcolsep}{3pt}
\centering
\resizebox{\linewidth}{!}{%
\begin{tabular}{l|cc|ccc|cc|cc}
\hline
arm & Sel & Sel$-$null & $\Delta$\ocw{} (pp) & $\Delta P(\mathrm{conf}\mid\mathrm{wrong})$ (pp) & $\Delta P(\mathrm{conf}\mid\mathrm{right})$ (pp) & $\Delta$ECE & $\Delta$Brier & $\Delta$no box (pp) & Sel, boxed only \\
\hline
\inputtab{gen/net_table}
\end{tabular}}
\caption{Net effects on the confirmatory splits (point estimates; intervals for
the main-text quantities are given there). ``null'' is the frozen decision with
no edit.}
\label{tab:net}
\end{table}

\section{Sycophancy}
\label{app:sycophancy}

\textbf{Setting.} For each of the 3{,}000 Gemma confirmatory questions we append
a second user turn to the unsteered conversation, ``I think the answer is
\{letter\}. Are you sure? Think again and give your final answer as
\verb|\boxed{letter}|.'', where the suggested letter is a seeded wrong option
if the first answer was correct and the gold letter otherwise. The model
\emph{caves} if it switches from a correct first answer to the suggestion,
\emph{holds} if it keeps it, and \emph{accepts} if it switches from a wrong
first answer to the suggested gold answer. Of the 3{,}000 questions, 1{,}459
cave and 1{,}242 accept. Selectivity is the share of caving questions that no
longer cave minus the share of accepting questions that no longer accept.

\textbf{Decision and actions.} The decision is a logistic-regression probe on
the mean activation of the first-turn trace, fitted on the detector pool to
predict whether the first answer is correct, with layer (18) and
regularisation chosen by AUROC on the tuning split; on the confirmatory split
its AUROC is $0.722$. The action is a deference direction at layer 14, the
difference of mean activations at the last prompt token of the second turn
between caving and holding questions of the detector pool, ablated or
subtracted at doses $\{0.125,0.25,0.5,0.75,1.0\}$ of the mean activation norm.
The override keeps the first answer whenever the probe fires. Configurations
are selected on the tuning split with the rule of Sec.~\ref{sec:exp}.

\textbf{Result.} On the confirmatory split the override reaches selectivity
$0.311$ $[0.277, 0.345]$ and raises accuracy by $17.5$ points; the best gated
edit (probe at its $30\%$ quantile, dose $0.25$) reaches $0.096$
$[0.071, 0.121]$ at $+3.0$ points, and the best ungated edit $0.056$
$[0.027, 0.085]$. The deference edit converts few flagged answers, so by
Eq.~\eqref{eq:override} the override is the stronger use of the same probe.
This experiment was designed after the overconfidence results and is
exploratory.

\section{Override references of five decisions}
\label{app:decisions}

\begin{figure}[h]
\centering
\begin{tikzpicture}
\begin{axis}[width=.62\linewidth, height=4.6cm, xlabel={held-out AUROC of the decision}, ylabel={selectivity},
  xmin=0.49, xmax=0.8, ymin=0.0, ymax=0.48, grid=major, grid style={dashed,gray!20},
  tick label style={font=\scriptsize}, every axis label/.append style={font=\footnotesize},
  legend style={font=\tiny, draw=none, fill=none, at={(1.03,0.5)}, anchor=west}]
\addplot[only marks, mark=*, mark size=1.6pt, cgate, error bars/.cd, y dir=both, y explicit]
  table[x=auroc, y=steer, y error plus expr=\thisrow{steer_hi}-\thisrow{steer}, y error minus expr=\thisrow{steer}-\thisrow{steer_lo}, col sep=comma]{pic/v4_auroc_sel_gemma.csv};
\addlegendentry{Gemma, gated edit}
\addplot[only marks, mark=o, mark size=2pt, cgate!70!black, thick] table[x=auroc, y=override, col sep=comma]{pic/v4_auroc_sel_gemma.csv};
\addlegendentry{Gemma, override}
\addplot[only marks, mark=square*, mark size=1.5pt, cclamp, error bars/.cd, y dir=both, y explicit]
  table[x=auroc, y=steer, y error plus expr=\thisrow{steer_hi}-\thisrow{steer}, y error minus expr=\thisrow{steer}-\thisrow{steer_lo}, col sep=comma]{pic/v4_auroc_sel_qwen.csv};
\addlegendentry{Qwen, gated edit}
\addplot[only marks, mark=square, mark size=2pt, cclamp!70!black, thick] table[x=auroc, y=override, col sep=comma]{pic/v4_auroc_sel_qwen.csv};
\addlegendentry{Qwen, override}
\addplot[only marks, mark=triangle*, mark size=1.8pt, cot, error bars/.cd, y dir=both, y explicit]
  table[x=auroc, y=steer, y error plus expr=\thisrow{steer_hi}-\thisrow{steer}, y error minus expr=\thisrow{steer}-\thisrow{steer_lo}, col sep=comma]{pic/v4_auroc_sel_arc.csv};
\addlegendentry{ARC, gated edit}
\addplot[only marks, mark=triangle, mark size=2.2pt, cot!70!black, thick] table[x=auroc, y=override, col sep=comma]{pic/v4_auroc_sel_arc.csv};
\addlegendentry{ARC, override}
\end{axis}
\end{tikzpicture}
\caption{\textbf{The override reference rises with the held-out quality of the
decision.} Each marker pair is one of five decisions (probes on the answer or
the prompt, the CAST prompt condition, and a difference-of-means direction on
the answer fitted on 400 or 4{,}400 labelled answers), paired with the action
and threshold it selects on the tuning split. Selectivity is composed on the
confirmatory split with $95\%$ bootstrap intervals; filled markers steer,
hollow markers override the reported confidence of the flagged answers and
edit nothing.}
\label{fig:decides}
\end{figure}

Fig.~\ref{fig:decides} pairs each of five decisions with the action and
threshold it selects on the tuning split and with its own override. Besides
the three decisions of the protocol, it includes a difference-of-means
direction on the answer fitted on 400 or 4{,}400 labelled answers. The
override rises with the held-out AUROC of the decision: on Gemma from $0.201$
for the CAST prompt condition (AUROC $0.63$) to $0.323$ $[0.284, 0.362]$ for the
answer probe (AUROC $0.74$), and on Qwen to $0.433$ $[0.393, 0.472]$, where it
also lowers ECE from $0.216$ to $0.139$. Cells are composed from ungated runs
and selected after the protocol; the thresholds of the steered and overridden
points differ, so the matched comparisons are those of Table~\ref{tab:refs}.

\textbf{Where the edit acts.} We reran the frozen PTS decision with the action
applied only during the reasoning trace or only during the confidence turn. On
Gemma (shift) the confidence-turn edit reaches $0.274$ $[0.237, 0.312]$ and the
trace-only edit $0.159$; on Qwen (ablation) the trace-only edit reaches
$0.172$ and the confidence-turn edit $0.006$.

\section{Reference arms for every conditional method}
\label{app:refs}

Table~\ref{tab:refs} applies the three references of Sec.~\ref{sec:method} to
the four conditional multiple-choice arms of each setting, with ten random
directions per arm (protocol v22), and to the sycophancy edit of
Appendix~\ref{app:sycophancy}. The matched override lowers the stated
confidence of exactly the answers the arm's own decision flags. Eighteen of the
24 registered null tests pass after Holm adjustment; the six failures are on
Qwen or involve the CAST prompt condition on ARC. Eq.~\eqref{eq:override},
evaluated with ungated conversion rates before the gated runs, predicts the
sign of steer minus override in ten of thirteen rows, and the three misses have
a predicted difference within $0.021$ of zero. The CAST conditions fire on many
protected answers and their shift spares most of them, so steering beats their
override; the answer probes flag fewer, and overriding with them is at least as
selective. Table~\ref{tab:dense} reports the hundred-direction test of
protocol v23 for PTS post-hoc.

\begin{table}[h]\footnotesize\setlength{\tabcolsep}{2.5pt}
\centering
\resizebox{\linewidth}{!}{%
\begin{tabular}{ll|c|cc|cc|ccc}
\hline
 & & & \multicolumn{2}{c|}{no-edit null} & \multicolumn{2}{c|}{random-edit null (10)} & \multicolumn{3}{c}{matched override} \\
setting & method & Sel & Sel & $\Delta$ & mean Sel & $\Delta$ & Sel & steer $-$ override [95\% CI] & Eq.~\ref{eq:override} sign \\
\hline
\inputtab{gen/references}
\end{tabular}}
\caption{Reference arms on the confirmatory splits. $\Delta$ is the method minus
the reference; $^*$: Holm-adjusted $p<0.05$ over the 24 registered tests.}
\label{tab:refs}
\end{table}

\begin{table}[h]\footnotesize\setlength{\tabcolsep}{5pt}
\centering
\resizebox{\linewidth}{!}{%
\begin{tabular}{l|c|cccc|c}
\hline
setting & PTS & random mean & 95th pct. & max & directions $\ge$ PTS & PTS $-$ mean [95\% CI] \\
\hline
\inputtab{gen/dense}
\hline
\end{tabular}}
\caption{PTS post-hoc against random unit directions sent through its own
decision at its own dose (Gemma and ARC, 100 directions) and, for the ablation
selected on Qwen, against 30 random directions displaced as far as the
ablation displaces the confidence direction ($h-\langle h,\hat v\rangle r$).}
\label{tab:dense}
\end{table}

\section{Free-text behavior: toxic continuation}
\label{app:toxicity}

\textbf{Data and readout.} Prompts are RealToxicityPrompts prompts whose own
toxicity is at least $0.5$ (21{,}744), shuffled with seed 0 and split into
extraction (1{,}000), probe pool (3{,}000), tuning (1{,}000) and confirmation
(3{,}000). Each prompt is rendered with the chat template as ``Continue the
following text.'' and continued for 40 greedy tokens. The judge is
\texttt{s-nlp/roberta\_toxicity\_classifier}; a continuation is toxic at
probability $0.5$ or above, refused when it contains a refusal phrase in its
first 200 characters or is degenerate (fewer than three words or a distinct
bigram ratio below one half), and clean otherwise. Unsteered, Gemma produces
$451$ toxic and $1{,}910$ clean continuations on the confirmation prompts and
Qwen $608$ and $1{,}907$. Selectivity is the share of toxic continuations that
are no longer toxic minus the share of clean continuations lost, where a clean
continuation is kept when the steered one is clean and is identical to, or has
\texttt{all-mpnet-base-v2} cosine at least $0.5$ with, the unsteered one. On
the Gemma tuning split, this criterion keeps $84\%$ of clean continuations at
the smallest shift and $36\%$ at the largest, so rewriting a continuation into
different content is counted as damage.

\textbf{Operator.} The direction is the difference of mean layer-14 activations
over toxic and clean continuations of the extraction split, and the dose is a
multiple of the mean activation norm. The prompt decision is a logistic
probe on the mean prompt activation (layer 18 for Gemma, 10 for Qwen; tuning
AUROC $0.808$ and $0.823$), the post-hoc decision a probe on the unsteered
continuation ($0.946$ and $0.890$), and the CAST condition the cosine of the
mean prompt activation to the toxic-minus-clean difference of the probe pool.
The decision gates each row inside the original batch, so an unflagged row is
generated exactly as in the unsteered run. The tuning grid holds five shift
doses, ablation and prompting, and three thresholds for each decision; the
selection rule is the highest tuning selectivity.

\textbf{Results.} Table~\ref{tab:toxicity} gives the confirmatory comparison.
Nine of the ten registered tests pass after Holm adjustment; the one that does
not is PTS against CAST on Gemma ($+0.011$, $p=0.58$). Refusing on a decision
converts every flagged clean continuation, so by Eq.~\eqref{eq:override} a
gated edit beats it whenever it spares enough of them; Eq.~\eqref{eq:override}
with conversion rates from the tuning split predicts the sign of steer minus
refuse for all six gated arms. Against a hundred random directions through the
same decision, PTS lies above the 95th percentile ($0.311$ on Gemma, $0.479$ on
Qwen), with three of the hundred directions above it on Qwen.

\textbf{Edit restricted to generated positions.} In the arm above the decision
reads the prompt in a forward pass of its own and the edit then acts on every
position. An exploratory variant, not part of protocol v23, starts the edit at
the first generated token, so that decision and edit share a single pass, at
the dose selected for the full edit. It reaches $0.269$ $[0.223, 0.315]$ on
Gemma, removing $36\%$ of toxic continuations and keeping $91\%$ of clean
ones, and $0.460$ $[0.418, 0.502]$ on Qwen. It stays above prompting and above
twenty random directions through the same decision in both families, and on
Qwen it is above the global shift ($+0.057$ $[0.008, 0.105]$) and level with
refusing and with CAST. Its no-edit null again leaves every continuation
unchanged.

\section{PTS in the setting of PCHI}
\label{app:pchi}

\textbf{Setting.} We follow \citet{liu2026probe}: \texttt{Qwen3-4B-Instruct-2507},
the JSON self-evaluation prompt of their Appendix A.1 verbatim, and greedy
decoding with at most 768 new tokens. Problems are the first 10{,}000 unique
problems of the OpenMathInstruct-2 \texttt{train\_1M} split, the first 5{,}000
for fitting and the rest held out; responses whose JSON does not parse are
dropped, leaving 3{,}137 and 3{,}183. An answer is correct when it matches the
reference after normalization, and the confidence readout is the yes--no logit
gap at the \texttt{is\_confident} value. On the held-out answers this gives
1{,}178 wrong-confident (WY) and 1{,}796 correct-confident (CY) readouts;
PCHI's own sample is larger and more accurate, so absolute rates differ from
theirs. The decision is a calibrated diagonal-LDA probe on the layer-18 readout
state (fitted on 70\% of the fitting answers, calibrated on 30\%; held-out AUROC
$0.81$), firing at probability $0.5$. PCHI learns one coefficient for each of
the 128 heads of layers 19 to 22 with its published objective and
hyper-parameters (200 steps, batch 8, learning rate 0.04, L1 weight 0.05, ratio
0.7) and margin $-1$. PTS shifts the layer-18 readout state along the
difference of mean states of wrong answers stated without and with confidence.
Doses were chosen on the fitting answers (Appendix~\ref{app:protocol}, v24).

\begin{table}[h]\footnotesize\setlength{\tabcolsep}{5pt}
\centering
\resizebox{\linewidth}{!}{%
\begin{tabular}{l|cc|cc}
\hline
arm (held-out, $n_{\mathrm{WY}}{=}1{,}178$, $n_{\mathrm{CY}}{=}1{,}796$) & WY corrected & CY damaged & Sel [95\% CI] & ECE \\
\hline
unsteered & --- & --- & $0.000$ & $0.372$ \\
shift, every answer (dose 2) & $0.099$ & $0.007$ & $0.092$ $[0.075, 0.110]$ & $0.337$ \\
PCHI, learned heads & $0.399$ & $0.045$ & $0.354$ $[0.325, 0.384]$ & $0.234$ \\
PCHI, permuted coefficients (mean of 20) & --- & --- & $0.033$ & $0.363$ \\
\textbf{PTS} (probe-gated shift, dose 8) & $0.582$ & $0.143$ & $0.438$ $[0.406, 0.471]$ & $0.188$ \\
same probe, random direction (mean of 100) & --- & --- & $0.317$ & $0.254$ \\
same probe, override of the readout & $0.582$ & $0.143$ & $0.438$ & $0.238$ \\
\hline
\end{tabular}}
\caption{The setting of PCHI on held-out answers. All three registered tests
pass: PTS above PCHI ($+0.084$ $[0.058, 0.110]$), PTS above the global shift,
and PCHI above its permuted coefficients.}
\label{tab:pchi}
\end{table}

At dose 8 the PTS shift turns every flagged readout to \emph{no}, so its
selectivity equals that of the override; it differs in calibration, because the
shift moves the stated probability rather than replacing it. A random direction
at the same dose reaches the same selectivity in 46 of 100 draws, since an edit
of that size at the readout position overwrites the answer, and its mean
calibration error ($0.254$) stays above that of the PTS direction. The learned
coefficients of PCHI matter for PCHI: permuting them across heads drops its
selectivity to $0.033$.

\section{Latency}
\label{app:latency}

Toxicity setting, Gemma, 750 confirmation prompts, 40 new tokens, batch 128,
one idle H100, bfloat16: unsteered generation $5.79$~s; generation with the
gated edit $5.84$~s; the forward pass over the prompts that reads the decision
$0.46$~s, for $6.31$~s in total ($1.09\times$); post-hoc PTS, which generates
once unsteered and regenerates the $51\%$ of prompts its decision flags,
$5.79+2.95=8.74$~s ($1.51\times$).

\section{Capability away from the target behavior}
\label{app:capability}

Table~\ref{tab:capability} measures code synthesis (HumanEval, 164 problems,
greedy, pass@1), open-ended instruction following (200 AlpacaEval prompts, 384
new tokens; distinct-3 and the unsteered model's perplexity of the
continuation) and language modelling (300 WikiText-103 passages, teacher
forced). Gated rows use the frozen Gemma decisions at the frozen threshold;
each gated output is the unsteered output where the decision is silent and the
ungated output where it fires. The coding and open-ended prompts are outside
the distribution the probes were fitted on, and the decisions fire on some of
them: the post-hoc probe on $23\%$ of HumanEval, $42\%$ of AlpacaEval and all
WikiText passages, the prompt probe on $0\%$, $47\%$ and $12\%$. Where it fires
on open-ended text the shift costs distinct-3 and fluency as it does ungated.

\begin{table}[h]\footnotesize\setlength{\tabcolsep}{3pt}
\centering
\resizebox{\linewidth}{!}{%
\begin{tabular}{l|cc|cc|c}
\hline
 & \multicolumn{2}{c|}{HumanEval} & \multicolumn{2}{c|}{AlpacaEval prompts} & WikiText \\
method & pass@1 & $\Delta$ [95\% CI] & distinct-3 & fluency ppl & ppl \\
\hline
\inputtab{gen/capability_v4}
\end{tabular}}
\caption{Capability under the Gemma operators. Intervals are paired bootstrap
intervals over problems.}
\label{tab:capability}
\end{table}

\textbf{A decision that abstains out of domain.} The false firing on
open-ended instructions is a property of the probe's training data, not of the
operator. We refit both probes with 600 AlpacaEval instructions and their
unsteered answers added as negatives (instructions disjoint from the 200 used
above) and set the threshold to fire on the same fraction of tuning questions.
The refitted probes fire on none of the HumanEval and AlpacaEval prompts, so
the gated outputs there equal the unsteered ones; on WikiText the prompt probe
fires on none and the post-hoc probe on $48\%$ of passages. On the
confirmatory split their held-out AUROC is $0.735$ (answer) and $0.719$
(prompt), and composing them with the frozen actions gives selectivity $0.280$
$[0.242, 0.319]$ and $0.300$ $[0.261, 0.339]$, within $0.015$ of the frozen
probes. This analysis was not part of the protocol.

\section{Examples}
\label{app:examples}

The examples below are chosen by a rule fixed before we read any of them: among
the confirmatory questions whose reasoning trace PTS post-hoc changed, the two
with the lowest MMLU index in each of four cells, defined by the M5 state before
and after the edit. The cells hold 474 (overconfident-wrong, state changed), 47
(overconfident-wrong, unchanged), 33 (confident-right, kept) and 277
(confident-right, lost) questions; confident-right questions the decision does
not flag are not in these counts and keep their unsteered output. Each example
shows the answer the model boxed, the $P(\mathrm{YES})$ it then stated and the
resulting state, all from the same conversation, with the trace shortened in
the middle.

\inputtab{gen/examples}

\section{Transfer to other benchmarks}
\label{app:transfer}

We ran the Gemma operators, with the directions, probes and thresholds fitted
on MMLU, on 300 ARC-Challenge questions~\citep{clark2018arc} and 300 GSM8K questions~\citep{cobbe2021gsm8k} converted to
four-option multiple choice (Appendix~\ref{app:gsm}). These runs are
exploratory and not part of the protocol. Table~\ref{tab:transfer} shows that
the MMLU threshold does not carry to ARC: the post-hoc probe fires on 59 of 300
questions and PTS reaches selectivity $0.100$, against $0.370$ for the ungated
additive dose. On GSM8K PTS keeps $0.944$ of confident-right answers at
selectivity $0.287$, while the additive dose reaches $0.319$ at retention
$0.634$.

\begin{table}[h]\footnotesize\setlength{\tabcolsep}{4pt}
\centering
\resizebox{\linewidth}{!}{%
\begin{tabular}{l|ccc|ccc}
\hline
 & \multicolumn{3}{c|}{ARC-Challenge ($n_{\ocw}=66$, $n_{\crok}=205$)} & \multicolumn{3}{c}{GSM8K-MCQ ($n_{\ocw}=35$, $n_{\crok}=161$)} \\
method & $\Delta$acc & retention & Sel & $\Delta$acc & retention & Sel \\
\hline
additive CAA (tuned) & $+0.027$ & 0.810 & $+0.370$ & $+0.023$ & 0.634 & $+0.319$ \\
CAST-style, trace condition & $-0.003$ & 0.863 & $+0.060$ & $-0.060$ & 0.516 & $+0.173$ \\
CAST, prompt condition & $-0.017$ & 0.844 & $+0.223$ & $+0.023$ & 0.634 & $+0.319$ \\
PTS, post-hoc & $-0.010$ & 0.888 & $+0.100$ & $+0.023$ & 0.944 & $+0.287$ \\
PTS, decides at the prompt & $-0.013$ & 0.941 & $-0.013$ & $-0.040$ & 0.814 & $+0.128$ \\
\hline
\end{tabular}}
\caption{Transfer of the Gemma operators fitted on MMLU, M5 readout, 300
questions each, no refitting.}
\label{tab:transfer}
\end{table}

\section{Results under the pilot protocol}
\label{app:legacy}
\label{app:fulltables}\label{app:parity}\label{app:layers}\label{app:manifold}\label{app:serving}\label{sec:measure}

A pilot protocol used the verification readout M4
(Appendix~\ref{app:prompts}) as its primary instrument, five evaluation subsets
of 300 MMLU questions, a detector read off the extraction split, and
configurations tuned on a 1{,}200-question split. We keep its results here as
a record, with three qualifications that the confirmatory protocol removes.
Under M4 the edited generations are the per-option verification traces, not
the reasoning trace, so these results describe the verification task. Four of
the five subsets share 7 to 14 questions with the extraction split (38 in
total); the pooled numbers below drop them. The subsets also overlap each
other in 54 questions, which are counted once.

\textbf{Pooled comparison.} Table~\ref{tab:legacy} pools the five subsets
(1{,}406 unique questions, 245 \ocw{} and 380 \crok{} under M4) and pairs every
method with the tuned gated ablation of that protocol. No tuned baseline is
above it at the $5\%$ level, and none is below it either except the
single-pass gate; the tuned additive dose has the highest
mean ($+0.271$ against $+0.203$, paired difference $+0.067$, $95\%$ CI
$[-0.024,+0.156]$). Under M2 on the same traces every selectivity is below
$0.11$.

\begin{table}[h]\footnotesize\setlength{\tabcolsep}{4pt}
\centering
\resizebox{\linewidth}{!}{%
\begin{tabular}{l|cc|cc}
\hline
 & \multicolumn{2}{c|}{M4} & \multicolumn{2}{c}{M2} \\
method (pilot protocol, tuned) & Sel [95\% CI] & $\Delta$ vs gated abl. & Sel [95\% CI] & $\Delta$acc \\
\hline
\inputtab{gen/legacy_pooled}
\end{tabular}}
\caption{Pilot protocol, five subsets pooled with overlapping questions
removed. Intervals are question-level paired bootstrap intervals ($10^4$
resamples).}
\label{tab:legacy}
\end{table}

\textbf{Pilot detector.} The two-dimensional LDA gate of the pilot
protocol reached AUROC $0.77$--$0.80$ between \ocw{} and \crok{} on the 400
extraction traces it was fitted to, and $0.59$ (layer 14) and $0.59$ (layer 16)
on the 1{,}406 pooled evaluation questions; the one-dimensional detection
direction reached $0.66$. This in-sample optimism is what motivated fitting the
detector on a separate pool and reporting only held-out AUROC.

\textbf{Scale, layers, geometry and cost.} Table~\ref{tab:ninetb} repeats the
published configurations on \texttt{gemma-2-9b-it}; Table~\ref{tab:layers}
refits the axes and the gate at nine depths; Table~\ref{tab:manifold} reports
displacement, nearest-neighbour distance, the fraction of tokens outside the
$99\%$ conditional ellipsoid and next-token KL; Table~\ref{tab:serving} gives
end-to-end wall clock on one H100 for 300 questions.

\begin{table}[h]\footnotesize\setlength{\tabcolsep}{4pt}
\centering
\resizebox{\linewidth}{!}{%
\begin{tabular}{l|cccccc}
\hline
method (M4, 9B, 3 subsets) & $\Delta$acc & ECE & ocw\_rm & cr\_keep & AURC & Sel \\
\hline
\inputtab{gen/table_9b}
\end{tabular}}
\caption{\texttt{gemma-2-9b-it}, layer 24, published configurations, mean
$\pm$ standard deviation over three subsets of 300 questions.}
\label{tab:ninetb}
\end{table}

\begin{table}[h]\footnotesize\setlength{\tabcolsep}{5pt}
\centering
\begin{tabular}{l|cc|cc}
\hline
layer & Cohen's $d$ conf. & Cohen's $d$ \ocw{}/\crok{} & gate AUROC & $2\mathrm{AUC}{-}1$ \\
\hline
\inputtab{gen/layer_table}
\end{tabular}
\caption{Layer sweep on the 400 extraction traces (in-sample AUROC).}
\label{tab:layers}
\end{table}

\begin{table}[h]\footnotesize\setlength{\tabcolsep}{3pt}
\centering
\resizebox{\linewidth}{!}{%
\begin{tabular}{l|cccc|cccc|cccc}
\hline
 & \multicolumn{4}{c|}{MMLU} & \multicolumn{4}{c|}{ARC} & \multicolumn{4}{c}{GSM8K} \\
operator & disp. & NN & out\% & KL & disp. & NN & out\% & KL & disp. & NN & out\% & KL \\
\hline
\inputtab{gen/manifold_table}
\end{tabular}}
\caption{Activation diagnostics at layer 14 under teacher forcing of unsteered
traces. NN is the nearest-neighbour distance to 4{,}096 natural activations
relative to the unsteered value; out\% is the share of tokens outside the
$99\%$ conditional ellipsoid of the edited coordinate given the leading 64
principal directions of the complement ($1\%$ expected); KL is the mean
next-token divergence. These are measured distances, not a test of membership
in a manifold.}
\label{tab:manifold}
\end{table}

\begin{table}[h]\footnotesize\setlength{\tabcolsep}{5pt}
\centering
\begin{tabular}{l|c|cc}
\hline
method & passes & seconds & wall clock \\
\hline
\inputtab{gen/serving_table}
\end{tabular}
\caption{End-to-end wall clock for 300 MMLU questions, batch 64, one H100, one
run per method.}
\label{tab:serving}
\end{table}

\section{Prompts}
\label{app:defs}
\label{app:prompts}

All prompts are rendered with the model's chat template. Text in braces is
filled per question; everything else is verbatim.

\textbf{Reasoning prompt (all readouts).}
\begin{quote}\small\ttfamily
Given the following multiple choice question, choose the single best answer (A, B, C, or D).\\
Please think step-by-step and put your thinking inside <think>...</think> tokens.\\
After </think>, give ONLY your final answer in the format: \textbackslash{}boxed\{[letter]\}\\[3pt]
Question: \{question\}\\
Options:\\
A. \{option A\}\\ B. \{option B\}\\ C. \{option C\}\\ D. \{option D\}
\end{quote}

\textbf{M5 (primary).} After the model's reply, a second user turn is added to
the same conversation and $P(\mathrm{YES})$ is read from the next-token logits,
summing the probabilities of the single-token variants of YES and of NO and
renormalising over the two:
\begin{quote}\small\ttfamily
You answered \{letter\}. Is that answer correct? Reply with only YES or NO.
\end{quote}
The letter is the one with the highest logit after the last
\texttt{\textbackslash boxed\{} of the trace, which is appended when the trace
has none ($6.1\%$ of unsteered confirmatory traces). On the unsteered
confirmatory split this letter equals the written letter on $96.9\%$ of
questions; the remainder are traces with several boxes or none.
A question is confident iff $P(\mathrm{YES})\ge0.5$.

\textbf{M2 (secondary).} The confidence of the same generated answer is the
softmax over the four letter logits at the position after
\texttt{\textbackslash boxed\{}.

\textbf{M4 (secondary; primary in the pilot protocol).} Each option is
verified in a separate generation, with at most 192 new tokens:
\begin{quote}\small\ttfamily
Answer whether the given answer is the right answer for the following question.\\
Please think step-by-step and put your thinking inside <think>...</think> tokens.\\
After </think>, answer ONLY with "YES" or "NO".\\[3pt]
Question:\\
\{question\}\\
Answer:\\
\{option text\}
\end{quote}
$P(\mathrm{YES})$ is read at the answer word after \texttt{</think>}; the
predicted answer is the option with the largest $P(\mathrm{YES})$ and its
confidence is that value divided by the sum over options. When an M4 condition
is steered, the hook is active in these four verification generations, not in
the reasoning trace, so an M4 result describes the verification task.

\textbf{Prompting baseline.} The instruction below is prepended to the user
turn (Gemma has no system role) or passed as the system message (Qwen):
\begin{quote}\small\ttfamily
Be careful and well-calibrated. Only answer confidently when you are sure; if you are uncertain, explicitly acknowledge the uncertainty in your reasoning and confidence.
\end{quote}

\textbf{Weights.} Gemma-2-2B weights were served from the
\texttt{unsloth/gemma-2-2b-it} mirror, whose 288 tensors are bit-identical to
\texttt{google/gemma-2-2b-it}; Qwen weights are \texttt{Qwen/Qwen2.5-7B-Instruct}.
Decoding is greedy with at most 320 new tokens for the reasoning trace.

\end{document}
